% Options for packages loaded elsewhere
\PassOptionsToPackage{unicode}{hyperref}
\PassOptionsToPackage{hyphens}{url}
\PassOptionsToPackage{dvipsnames,svgnames,x11names}{xcolor}
\documentclass[
  italian,
  10pt,
]{book}
\usepackage{xcolor}
\usepackage[papersize={6in,9in},inner=0.75in,outer=0.5in,top=0.7in,bottom=0.7in]{geometry}
\usepackage{amsmath,amssymb}
\setcounter{secnumdepth}{-\maxdimen} % remove section numbering
\usepackage{iftex}
\ifPDFTeX
  \usepackage[T1]{fontenc}
  \usepackage[utf8]{inputenc}
  \usepackage{textcomp} % provide euro and other symbols
\else % if luatex or xetex
  \usepackage{unicode-math} % this also loads fontspec
  \defaultfontfeatures{Scale=MatchLowercase}
  \defaultfontfeatures[\rmfamily]{Ligatures=TeX,Scale=1}
\fi
\usepackage{lmodern}
\ifPDFTeX\else
  % xetex/luatex font selection
\fi
% Use upquote if available, for straight quotes in verbatim environments
\IfFileExists{upquote.sty}{\usepackage{upquote}}{}
\IfFileExists{microtype.sty}{% use microtype if available
  \usepackage[]{microtype}
  \UseMicrotypeSet[protrusion]{basicmath} % disable protrusion for tt fonts
}{}
\makeatletter
\@ifundefined{KOMAClassName}{% if non-KOMA class
  \IfFileExists{parskip.sty}{%
    \usepackage{parskip}
  }{% else
    \setlength{\parindent}{0pt}
    \setlength{\parskip}{6pt plus 2pt minus 1pt}}
}{% if KOMA class
  \KOMAoptions{parskip=half}}
\makeatother
\usepackage{color}
\usepackage{fancyvrb}
\newcommand{\VerbBar}{|}
\newcommand{\VERB}{\Verb[commandchars=\\\{\}]}
\DefineVerbatimEnvironment{Highlighting}{Verbatim}{commandchars=\\\{\}}
% Add ',fontsize=\small' for more characters per line
\usepackage{framed}
\definecolor{shadecolor}{RGB}{248,248,248}
\newenvironment{Shaded}{\begin{snugshade}}{\end{snugshade}}
\newcommand{\AlertTok}[1]{\textcolor[rgb]{0.94,0.16,0.16}{#1}}
\newcommand{\AnnotationTok}[1]{\textcolor[rgb]{0.56,0.35,0.01}{\textbf{\textit{#1}}}}
\newcommand{\AttributeTok}[1]{\textcolor[rgb]{0.13,0.29,0.53}{#1}}
\newcommand{\BaseNTok}[1]{\textcolor[rgb]{0.00,0.00,0.81}{#1}}
\newcommand{\BuiltInTok}[1]{#1}
\newcommand{\CharTok}[1]{\textcolor[rgb]{0.31,0.60,0.02}{#1}}
\newcommand{\CommentTok}[1]{\textcolor[rgb]{0.56,0.35,0.01}{\textit{#1}}}
\newcommand{\CommentVarTok}[1]{\textcolor[rgb]{0.56,0.35,0.01}{\textbf{\textit{#1}}}}
\newcommand{\ConstantTok}[1]{\textcolor[rgb]{0.56,0.35,0.01}{#1}}
\newcommand{\ControlFlowTok}[1]{\textcolor[rgb]{0.13,0.29,0.53}{\textbf{#1}}}
\newcommand{\DataTypeTok}[1]{\textcolor[rgb]{0.13,0.29,0.53}{#1}}
\newcommand{\DecValTok}[1]{\textcolor[rgb]{0.00,0.00,0.81}{#1}}
\newcommand{\DocumentationTok}[1]{\textcolor[rgb]{0.56,0.35,0.01}{\textbf{\textit{#1}}}}
\newcommand{\ErrorTok}[1]{\textcolor[rgb]{0.64,0.00,0.00}{\textbf{#1}}}
\newcommand{\ExtensionTok}[1]{#1}
\newcommand{\FloatTok}[1]{\textcolor[rgb]{0.00,0.00,0.81}{#1}}
\newcommand{\FunctionTok}[1]{\textcolor[rgb]{0.13,0.29,0.53}{\textbf{#1}}}
\newcommand{\ImportTok}[1]{#1}
\newcommand{\InformationTok}[1]{\textcolor[rgb]{0.56,0.35,0.01}{\textbf{\textit{#1}}}}
\newcommand{\KeywordTok}[1]{\textcolor[rgb]{0.13,0.29,0.53}{\textbf{#1}}}
\newcommand{\NormalTok}[1]{#1}
\newcommand{\OperatorTok}[1]{\textcolor[rgb]{0.81,0.36,0.00}{\textbf{#1}}}
\newcommand{\OtherTok}[1]{\textcolor[rgb]{0.56,0.35,0.01}{#1}}
\newcommand{\PreprocessorTok}[1]{\textcolor[rgb]{0.56,0.35,0.01}{\textit{#1}}}
\newcommand{\RegionMarkerTok}[1]{#1}
\newcommand{\SpecialCharTok}[1]{\textcolor[rgb]{0.81,0.36,0.00}{\textbf{#1}}}
\newcommand{\SpecialStringTok}[1]{\textcolor[rgb]{0.31,0.60,0.02}{#1}}
\newcommand{\StringTok}[1]{\textcolor[rgb]{0.31,0.60,0.02}{#1}}
\newcommand{\VariableTok}[1]{\textcolor[rgb]{0.00,0.00,0.00}{#1}}
\newcommand{\VerbatimStringTok}[1]{\textcolor[rgb]{0.31,0.60,0.02}{#1}}
\newcommand{\WarningTok}[1]{\textcolor[rgb]{0.56,0.35,0.01}{\textbf{\textit{#1}}}}
\usepackage{longtable,booktabs,array}
\usepackage{caption}
\captionsetup[table]{skip=6pt}
\newcounter{none} % for unnumbered tables
\usepackage{calc} % for calculating minipage widths
% Correct order of tables after \paragraph or \subparagraph
\usepackage{etoolbox}
\makeatletter
\patchcmd\longtable{\par}{\if@noskipsec\mbox{}\fi\par}{}{}
\makeatother
\ifLuaTeX
\usepackage[bidi=basic,shorthands=off]{babel}
\else
\usepackage[bidi=default,shorthands=off]{babel}
\fi
\ifLuaTeX
  \usepackage{selnolig} % disable illegal ligatures
\fi
\setlength{\emergencystretch}{3em} % prevent overfull lines
\providecommand{\tightlist}{%
  \setlength{\itemsep}{0pt}\setlength{\parskip}{0pt}}
\usepackage{fvextra}
\usepackage{float}
\usepackage{xurl}
\usepackage[htt]{hyphenat}
\DefineVerbatimEnvironment{Highlighting}{Verbatim}{breaklines,breakanywhere,fontsize=\small,commandchars=\\\{\}}
\floatplacement{figure}{H}
\setlength{\emergencystretch}{8em}
\sloppy
\usepackage{bookmark}
\IfFileExists{xurl.sty}{\usepackage{xurl}}{} % add URL line breaks if available
\urlstyle{same}
% fallback for those not using the hyperref driver hyperxmp:
\makeatletter
\@ifundefined{xmpquote}{\newcommand{\xmpquote}[1]{#1}}{}
\makeatother
\hypersetup{
  pdftitle={Quaderno AI},
  pdfauthor={Andrea Merlin},
  pdflang={it},
  colorlinks=true,
  linkcolor={blue},
  filecolor={Maroon},
  citecolor={Blue},
  urlcolor={blue},
  pdfcreator={LaTeX via pandoc}}

\title{Quaderno AI}
\usepackage{etoolbox}
\makeatletter
\providecommand{\subtitle}[1]{% add subtitle to \maketitle
  \apptocmd{\@title}{\par {\large #1 \par}}{}{}
}
\makeatother
\subtitle{Concetti, strumenti e percorsi pratici per lo sviluppo
software con AI}
\author{Andrea Merlin}
\date{Numero 1 - Novembre 2026}

\begin{document}
\frontmatter
\maketitle

{
\setcounter{tocdepth}{1}
\tableofcontents
}
\mainmatter
\chapter{Capitolo 5}\label{capitolo-5}

\section{Migrare a .NET da altri linguaggi e
piattaforme}\label{migrare-a-.net-da-altri-linguaggi-e-piattaforme}

\textbf{Abstract.} Migrare un'applicazione non significa riscrivere ogni
riga cambiando sintassi. Significa trasferire in un nuovo ecosistema un
insieme di comportamenti, dati, contratti, vincoli operativi e
aspettative degli utenti, riducendo il rischio di interrompere il
servizio. Questo capitolo guida chi conosce già PHP, Python,
JavaScript/Node.js, Java o un altro stack attraverso il passaggio a .NET
moderno e ASP.NET Core. Partiremo dall'assessment del sistema esistente
e confronteremo le strategie di migrazione; poi costruiremo una mappa
dei concetti, esploreremo hosting, pipeline HTTP, Dependency Injection,
dati, sicurezza, test, prestazioni e deployment. Un caso pratico
mostrerà come migrare un servizio per passi, mantenendo compatibili i
client e predisponendo rollback e osservabilità. L'obiettivo non è
dichiarare un framework universalmente migliore, ma fornire un metodo
verificabile per scegliere cosa trasferire, cosa adattare e cosa
lasciare temporaneamente nel sistema originale.

\subsection{Cosa imparerai in questo
capitolo}\label{cosa-imparerai-in-questo-capitolo}

Al termine saprai:

\begin{itemize}
\tightlist
\item
  distinguere una migrazione tecnologica da una riscrittura funzionale e
  raccogliere le informazioni necessarie prima di stimare il lavoro;
\item
  confrontare PHP, Python, Node.js e Java con C\# e .NET senza
  confondere analogie utili con equivalenze perfette;
\item
  scegliere tra riscrittura completa, migrazione per moduli e
  sostituzione incrementale con il pattern \emph{strangler};
\item
  creare e organizzare un servizio ASP.NET Core, capire il ruolo
  dell'host, di Kestrel e di \texttt{Program.cs};
\item
  trasferire route, middleware, validazione, error handling,
  configurazione e stato;
\item
  pianificare schema, migrazioni e dati con Entity Framework Core
  mantenendo le compatibilità necessarie;
\item
  migrare autenticazione, autorizzazioni e segreti senza abbassare il
  livello di protezione;
\item
  costruire test che descrivano il comportamento attuale e verificare la
  parità funzionale;
\item
  misurare prestazioni, integrare CI/CD, distribuire in container e
  preparare rollout e rollback;
\item
  usare checklist e criteri di uscita per dimostrare che la migrazione è
  pronta.
\end{itemize}

\subsection{Prerequisiti}\label{prerequisiti}

\begin{itemize}
\tightlist
\item
  Conoscenza di base di almeno un linguaggio e framework di sviluppo.
\item
  Familiarità con HTTP, API, database relazionali, test automatici e
  controllo versione.
\item
  Un SDK .NET 10 installato per eseguire gli esempi. Non è richiesta
  esperienza precedente con C\# o ASP.NET Core.
\end{itemize}

\subsection{Convenzioni}\label{convenzioni}

\begin{itemize}
\tightlist
\item
  I comandi sono mostrati in forma multipiattaforma con \texttt{dotnet};
  i frammenti di shell possono richiedere piccoli adattamenti in
  PowerShell.
\item
  Il riferimento è .NET 10, salvo indicazione diversa. Le API e i
  template possono variare tra versioni: verificare sempre il target
  dichiarato nel \texttt{.csproj}.
\item
  Gli esempi sono didattici: prima di usarli in produzione occorre
  adattare autenticazione, segreti, storage, log e policy operative.
\item
  ``Sistema sorgente'' indica l'applicazione da cui si migra; ``sistema
  target'' indica quella realizzata con .NET. ``Cutover'' indica il
  passaggio del traffico al target.
\end{itemize}

\subsection{Indice concettuale}\label{indice-concettuale}

\begin{enumerate}
\def\labelenumi{\arabic{enumi}.}
\tightlist
\item
  Perché migrare e come valutare il progetto
\item
  Mappare linguaggi, runtime e strumenti
\item
  Pianificare una migrazione controllata
\item
  Architettura e organizzazione del codice
\item
  Richieste HTTP e pipeline
\item
  Configurazione, stato e dipendenze
\item
  Dati e integrazione con i sistemi esistenti
\item
  Sicurezza e identità
\item
  Testing e qualità
\item
  Prestazioni e comportamento del runtime
\item
  Container, CI/CD e produzione
\item
  Caso pratico: migrare un servizio catalogo
\item
  Approfondire la mappa dai principali ecosistemi
\item
  Prove di migrazione dei dati: runbook e controlli
\item
  Diagnosi operativa durante la migrazione
\item
  Verifica finale, criteri di successo e fonti ufficiali
\item
  Checklist, glossario e anti-pattern
\item
  Appendice: piano di migrazione riutilizzabile
\end{enumerate}

\section{5.1 Perché migrare e come valutare il
progetto}\label{perchuxe9-migrare-e-come-valutare-il-progetto}

Una migrazione viene spesso proposta come risposta a un problema
percepito: un framework è fuori supporto, il team non riesce più a
trovare competenze, il costo di manutenzione cresce o il servizio deve
entrare in un ambiente standardizzato su .NET. La proposta tecnica,
però, non dimostra da sola che la migrazione sia vantaggiosa. Prima di
scegliere la destinazione bisogna comprendere il sistema che si intende
sostituire, il valore che offre e i vincoli che non sono visibili
osservando soltanto i file sorgente.

\subsection{5.1.1 Migrazione non è
traduzione}\label{migrazione-non-uxe8-traduzione}

Un'applicazione è più del suo codice. Comprende almeno:

\begin{itemize}
\tightlist
\item
  comportamenti osservabili, inclusi casi limite e comportamenti
  accidentali su cui i client hanno finito per fare affidamento;
\item
  contratti HTTP, formati, codici di stato, eventi, code, file e schemi
  dati;
\item
  identità, ruoli, permessi, credenziali e regole di conservazione;
\item
  operazioni pianificate, procedure manuali e strumenti di supporto;
\item
  caratteristiche operative: latenza, disponibilità, throughput, costo e
  procedure di ripristino.
\end{itemize}

Trasferire ogni funzione una a una può perpetuare difetti e vincoli non
più necessari. Cambiare tutto insieme, d'altra parte, rende difficile
capire se un malfunzionamento dipende dalla nuova piattaforma o da una
modifica funzionale. In una migrazione tecnologica ben controllata si
definisce prima quali comportamenti devono restare invariati. Le nuove
funzionalità e le correzioni di dominio possono essere introdotte
separatamente, oppure identificate esplicitamente come variazioni di
contratto.

Una distinzione utile è tra \textbf{parità funzionale} e \textbf{parità
operativa}. La prima riguarda ciò che utenti e client possono fare; la
seconda riguarda come il servizio si comporta sotto carico, come viene
distribuito e osservato, quali procedure servono per ripristinarlo. Un
endpoint che restituisce lo stesso JSON ma non rispetta più il limite di
latenza del client è funzionalmente simile, non una sostituzione
riuscita.

\subsection{5.1.2 Motivazioni e
controindicazioni}\label{motivazioni-e-controindicazioni}

Tra le motivazioni ragionevoli ci sono il termine del supporto del
runtime, la difficoltà di aggiornare dipendenze critiche, i requisiti di
sicurezza o conformità, l'esigenza di unificare toolchain e
infrastruttura, l'integrazione con librerie .NET già adottate e la
necessità di sostenere un carico o una modalità di deployment diversi.

Non sono sufficienti, da sole, frasi come ``C\# è più veloce'', ``il
linguaggio è vecchio'' o ``vogliamo un'architettura moderna''.
Un'applicazione stabile, ben testata e con costo operativo contenuto può
essere più sicura da aggiornare sullo stack corrente. Una riscrittura
consuma capacità di sviluppo, crea un periodo di doppia manutenzione e
può introdurre regressioni. Il confronto deve includere il costo totale:
sviluppo, licenze, hosting, supporto, formazione, incidenti e
dismissione del vecchio sistema.

La domanda corretta non è «.NET è migliore?», ma «quali requisiti del
nostro sistema soddisfa meglio il target, e con quali costi e rischi
misurabili?». Un prototipo può validare un'integrazione o una libreria;
non prova automaticamente che una migrazione completa sia conveniente.

\subsection{5.1.3 Inventario tecnico e
funzionale}\label{inventario-tecnico-e-funzionale}

L'assessment inizia da un inventario verificabile. Non basta contare
controller e righe di codice. Raccogliere, dove disponibili:

\begin{enumerate}
\def\labelenumi{\arabic{enumi}.}
\tightlist
\item
  \textbf{Superfici funzionali:} route, metodi, parametri, payload,
  errori, pagine e client che le usano.
\item
  \textbf{Processi:} job schedulati, consumer di code, import/export,
  webhook, strumenti amministrativi e script manuali.
\item
  \textbf{Dati:} database, tabelle, volumi, crescita, vincoli, indici,
  backup, retention e migrazioni pregresse.
\item
  \textbf{Dipendenze:} runtime, framework, pacchetti, servizi cloud,
  SDK, sistemi di identità e componenti nativi.
\item
  \textbf{Integrazioni:} API esterne, code, SMTP, storage di file,
  pagamenti, sistemi legacy e formati di scambio.
\item
  \textbf{Operazioni:} pipeline, deployment, secret store, logging,
  monitoraggio, alert, runbook, RTO e RPO.
\item
  \textbf{Utilizzo:} frequenza delle route, profili di carico, picchi,
  client attivi, operazioni costose e stagionalità.
\item
  \textbf{Vincoli:} finestre di fermo, requisiti legali, localizzazione
  dei dati, compatibilità, competenze e budget.
\end{enumerate}

Accoppiare l'inventario statico con osservazioni runtime. Log, tracing e
metriche rivelano percorsi davvero usati e dipendenze effettive;
repository e documentazione rivelano intenzioni e configurazioni.
Nessuna delle due fonti è completa. Un route presente nel codice può
essere inutilizzato, ma non va rimosso senza verifica con i proprietari
dei client.

\subsection{5.1.4 Requisiti non funzionali e
baseline}\label{requisiti-non-funzionali-e-baseline}

Per valutare il target servono misure di partenza. Registrare almeno:

\begin{itemize}
\tightlist
\item
  latenza per operazioni rappresentative, includendo percentili come
  p50, p95 e p99;
\item
  richieste al secondo, concorrenza, code e frequenza degli errori;
\item
  consumo di CPU, memoria, storage, connessioni e risorse di database;
\item
  tempi di avvio, deploy, ripristino e build;
\item
  disponibilità, incidenti, obiettivi RTO/RPO e carico operativo del
  team.
\end{itemize}

Misurare le stesse operazioni con input e ambiente confrontabili. Una
media globale nasconde le route lente e la variabilità; un benchmark su
una singola macchina non rappresenta una produzione con rete, database e
cache. Annotare configurazione, versione, dimensione del dataset e data
della misura. Senza una baseline il team non può distinguere un
miglioramento reale da una differenza di ambiente.

Definire criteri d'accettazione prima dell'implementazione. Per esempio:
tutti i flussi critici hanno test di caratterizzazione; il target
rispetta la compatibilità API; p95 non supera il budget concordato; il
rollback è stato provato; nessun segreto è contenuto nell'immagine. I
numeri devono provenire dal servizio e dai suoi utenti, non da soglie
arbitrarie.

\subsection{5.1.5 Registro delle decisioni e dei
rischi}\label{registro-delle-decisioni-e-dei-rischi}

Per ogni funzione si può registrare proprietario, frequenza, criticità,
dipendenze e strategia proposta. Un registro decisionale breve dovrebbe
descrivere la decisione, le alternative considerate, le evidenze e le
condizioni che la farebbero rivedere. Esempio: mantenere inizialmente il
database invariato perché più servizi condividono lo schema;
riconsiderare la separazione solo dopo aver stabilizzato API e
ownership.

Valutare probabilità e impatto dei rischi: dati non interpretabili allo
stesso modo, dipendenza non portabile, identità non migrabili, traffico
non riproducibile, doppia scrittura incoerente, indisponibilità durante
il cutover o mancanza di rollback. Ogni rischio importante deve avere
una mitigazione e un responsabile. Se manca un'informazione critica,
trasformarla in un'attività di discovery con un risultato atteso, non in
una supposizione nascosta nel piano.

\textbf{Risultato dell'assessment:} una mappa del sistema, una baseline,
un insieme di contratti da preservare, una lista di incertezze e una
raccomandazione motivata tra migrare, modernizzare sullo stack esistente
o rimandare. Questo pacchetto è la base della pianificazione, non un
documento da archiviare senza seguito.

\section{5.2 Mappare linguaggi, runtime e
strumenti}\label{mappare-linguaggi-runtime-e-strumenti}

Chi arriva in .NET porta con sé un modello mentale valido per il proprio
ecosistema. Riutilizzarlo rende più rapido l'apprendimento, a patto di
non supporre che due strumenti con nomi o scopi simili abbiano lo stesso
lifecycle, la stessa semantica o gli stessi limiti. In questa sezione
costruiamo una mappa e una piccola applicazione che useremo nei capitoli
successivi.

\subsection{5.2.1 Da .NET Core a .NET
moderno}\label{da-.net-core-a-.net-moderno}

``.NET Core'' è il nome della linea open source multipiattaforma fino a
.NET Core 3.1. Da .NET 5 il prodotto si chiama \textbf{.NET}: .NET 6, 7,
8, 9 e 10 sono versioni della piattaforma unificata. \textbf{ASP.NET
Core} è il framework web moderno della piattaforma e ha mantenuto
``Core'' nel proprio nome. Il vecchio \textbf{.NET Framework} è un
prodotto distinto, principalmente legato a Windows; non è il target di
questo capitolo.

Per un nuovo servizio, fissare una versione supportata e dichiararla
esplicitamente. .NET 10 è il riferimento degli esempi. Non scegliere il
target perché il computer dello sviluppatore lo ha installato per caso:
SDK, target framework, pipeline CI e runtime di produzione devono essere
coordinati. Per la policy di supporto e le patch correnti fare
riferimento alla pagina ufficiale del ciclo di vita .NET.

\subsection{5.2.2 Concetti e strumenti a
confronto}\label{concetti-e-strumenti-a-confronto}

\begin{itemize}
\tightlist
\item
  \textbf{PHP:} linguaggio PHP; runtime PHP e Composer; progetto
  \texttt{composer.json}; avvio dipendente dal framework o dal worker;
  framework come Laravel e Symfony; pacchetti da Packagist.
\item
  \textbf{Python:} linguaggio Python; runtime e pip/Poetry; progetto
  \texttt{pyproject.toml}; avvio tramite modulo, ASGI o WSGI; framework
  come Django, Flask e FastAPI; pacchetti da PyPI.
\item
  \textbf{Node.js:} JavaScript o TypeScript; Node.js con npm/pnpm;
  progetto \texttt{package.json}; entry point Node; framework come
  Express e NestJS; pacchetti dal registry npm.
\item
  \textbf{Java:} linguaggio Java; JDK con Maven/Gradle; progetto
  \texttt{pom.xml} o file Gradle; avvio da \texttt{main} o framework;
  framework come Spring Boot; pacchetti da Maven Central.
\item
  \textbf{.NET:} linguaggio C\#; SDK e runtime .NET; progetto
  \texttt{.csproj}; avvio da \texttt{Program.cs} e host; framework come
  ASP.NET Core; pacchetti da NuGet.
\end{itemize}

La tabella orienta, ma non mappa capacità uno-a-uno. Per esempio un
provider del framework sorgente può combinare routing, dependency
injection, sessione e configurazione; nel target queste responsabilità
possono appartenere a servizi differenti. Prima di sostituirlo,
scomporre quali comportamenti effettivamente offre.

\subsection{5.2.3 SDK, runtime e CLI}\label{sdk-runtime-e-cli}

L'\textbf{SDK} include compilatore, template e strumenti per sviluppare;
il \textbf{runtime} esegue un'applicazione già pubblicata. Un server di
build richiede l'SDK; un container di produzione può usare solo il
runtime, oppure un artefatto self-contained a seconda della strategia.
La distinzione influenza immagini, dimensione del deploy, patching e
installazione.

La CLI standardizza molte operazioni:

\begin{Shaded}
\begin{Highlighting}[]
\NormalTok{dotnet {-}{-}info}
\NormalTok{dotnet new webapi {-}n Catalogo.Api}
\NormalTok{cd Catalogo.Api}
\NormalTok{dotnet restore}
\NormalTok{dotnet build}
\NormalTok{dotnet test}
\NormalTok{dotnet run}
\NormalTok{dotnet publish {-}c Release {-}o publish}
\end{Highlighting}
\end{Shaded}

\texttt{dotnet\ restore} risolve i pacchetti; \texttt{build} compila;
\texttt{test} esegue i test; \texttt{run} avvia il progetto;
\texttt{publish} produce l'output destinato alla distribuzione. Non sono
sinonimi: per esempio una build riuscita non garantisce che
l'applicazione sia configurata correttamente in produzione.

Per pinning riproducibile, il repository può contenere
\texttt{global.json} con la policy SDK usata dal team. Il file
\texttt{.csproj} dichiara il target:

\begin{Shaded}
\begin{Highlighting}[]
\NormalTok{\textless{}}\KeywordTok{Project} \OtherTok{Sdk=}\StringTok{"Microsoft.NET.Sdk.Web"}\NormalTok{\textgreater{}}
\NormalTok{  \textless{}}\KeywordTok{PropertyGroup}\NormalTok{\textgreater{}}
\NormalTok{    \textless{}}\KeywordTok{TargetFramework}\NormalTok{\textgreater{}net10.0\textless{}/}\KeywordTok{TargetFramework}\NormalTok{\textgreater{}}
\NormalTok{    \textless{}}\KeywordTok{Nullable}\NormalTok{\textgreater{}enable\textless{}/}\KeywordTok{Nullable}\NormalTok{\textgreater{}}
\NormalTok{    \textless{}}\KeywordTok{ImplicitUsings}\NormalTok{\textgreater{}enable\textless{}/}\KeywordTok{ImplicitUsings}\NormalTok{\textgreater{}}
\NormalTok{  \textless{}/}\KeywordTok{PropertyGroup}\NormalTok{\textgreater{}}
\NormalTok{\textless{}/}\KeywordTok{Project}\NormalTok{\textgreater{}}
\end{Highlighting}
\end{Shaded}

\texttt{TargetFramework} non è una lista di pacchetti. Le librerie
esterne si aggiungono con \texttt{dotnet\ add\ package} e vengono
registrate nel progetto; in una soluzione con più progetti si distingue
il progetto API dai progetti di dominio e test.

\subsection{5.2.4 Sintassi C\# per chi conosce già un
linguaggio}\label{sintassi-c-per-chi-conosce-giuxe0-un-linguaggio}

C\# è staticamente tipizzato e compilato. Il compilatore intercetta
categorie di errori prima dell'avvio, ma il sistema dei tipi non
sostituisce validazione e test. \texttt{string?} indica che un valore
può essere null quando nullable reference types sono attivi;
\texttt{string} esprime l'intenzione opposta. I dati ricevuti da HTTP
sono comunque input non attendibile e vanno validati.

Un record può descrivere un DTO immutabile:

\begin{Shaded}
\begin{Highlighting}[]
\KeywordTok{public} \KeywordTok{sealed}\NormalTok{ record }\FunctionTok{CreateProductRequest}\OperatorTok{(}\DataTypeTok{string}\NormalTok{ Name}\OperatorTok{,} \DataTypeTok{decimal}\NormalTok{ Price}\OperatorTok{);}
\end{Highlighting}
\end{Shaded}

La definizione non valida automaticamente \texttt{Name} o
\texttt{Price}: occorrono regole esplicite. Un metodo asincrono
restituisce normalmente \texttt{Task\textless{}T\textgreater{}} o
\texttt{ValueTask\textless{}T\textgreater{}} e usa \texttt{await};
l'asincronia è essenziale per non bloccare thread mentre si aspetta I/O,
ma non rende automaticamente veloce una query.

\begin{Shaded}
\begin{Highlighting}[]
\KeywordTok{public}\NormalTok{ async Task}\OperatorTok{\textless{}}\NormalTok{Product}\OperatorTok{?\textgreater{}} \FunctionTok{FindAsync}\OperatorTok{(}
    \DataTypeTok{int}\NormalTok{ id}\OperatorTok{,}\NormalTok{ CancellationToken cancellationToken}\OperatorTok{)}
\OperatorTok{\{}
    \KeywordTok{return}\NormalTok{ await db}\OperatorTok{.}\FunctionTok{Products}
        \OperatorTok{.}\FunctionTok{SingleOrDefaultAsync}\OperatorTok{(}\NormalTok{p }\OperatorTok{=\textgreater{}}\NormalTok{ p}\OperatorTok{.}\FunctionTok{Id} \OperatorTok{==}\NormalTok{ id}\OperatorTok{,}\NormalTok{ cancellationToken}\OperatorTok{);}
\OperatorTok{\}}
\end{Highlighting}
\end{Shaded}

Il codice non dovrebbe nascondere le eccezioni con
\texttt{catch\ \{\ return\ null;\ \}}: un errore database non equivale a
un prodotto assente. Gli errori devono essere tradotti in modo coerente
al confine applicativo e registrati con contesto sufficiente, senza
esporre dettagli interni al client.

\subsection{5.2.5 Lifecycle e hosting}\label{lifecycle-e-hosting}

Molti framework sorgente offrono un processo per richiesta, un worker
persistente, oppure più processi con un server applicativo davanti.
ASP.NET Core usa un \textbf{host} che costruisce configurazione,
logging, dependency injection e lifecycle dell'applicazione. Kestrel è
il server HTTP integrato; può ricevere direttamente traffico o lavorare
dietro un reverse proxy, in base all'ambiente.

Nel modello moderno \texttt{Program.cs} configura servizi e pipeline e
avvia l'host. La parte essenziale può essere molto breve:

\begin{Shaded}
\begin{Highlighting}[]
\DataTypeTok{var}\NormalTok{ builder }\OperatorTok{=}\NormalTok{ WebApplication}\OperatorTok{.}\FunctionTok{CreateBuilder}\OperatorTok{(}\NormalTok{args}\OperatorTok{);}

\DataTypeTok{var}\NormalTok{ app }\OperatorTok{=}\NormalTok{ builder}\OperatorTok{.}\FunctionTok{Build}\OperatorTok{();}
\NormalTok{app}\OperatorTok{.}\FunctionTok{MapGet}\OperatorTok{(}\StringTok{"/health"}\OperatorTok{,} \OperatorTok{()} \OperatorTok{=\textgreater{}}\NormalTok{ Results}\OperatorTok{.}\FunctionTok{Ok}\OperatorTok{(}\KeywordTok{new} \OperatorTok{\{}\NormalTok{ status }\OperatorTok{=} \StringTok{"ok"} \OperatorTok{\}));}
\NormalTok{app}\OperatorTok{.}\FunctionTok{Run}\OperatorTok{();}
\end{Highlighting}
\end{Shaded}

\texttt{AddEndpointsApiExplorer()} da solo registra metadati, ma non
genera né espone un documento OpenAPI: serve insieme a un generatore.
Per esempio, in .NET 10, aggiungendo il pacchetto
\texttt{Microsoft.AspNetCore.OpenApi}, si può usare
\texttt{AddOpenApi()} con \texttt{MapOpenApi()}; qui non è necessario e
viene omesso.

La brevità non implica assenza di struttura: route, servizi, opzioni e
componenti possono essere estratti in metodi o moduli. Il punto
d'ingresso non è il luogo in cui inserire tutta la logica di dominio; è
il punto in cui si compone l'applicazione.

\subsection{5.2.6 Pacchetti e dipendenze}\label{pacchetti-e-dipendenze}

NuGet è il sistema di pacchetti .NET. Ogni dipendenza aggiunta aumenta
superficie di aggiornamento, licenze e rischio di supply chain. Prima di
portare un pacchetto, verificare compatibilità con il target,
manutenzione, advisory, licenza, dipendenze transitive e necessità
effettiva. Un equivalente con nome simile può avere una filosofia o un
lifecycle differenti.

Per build riproducibili, versionare il file di lock quando il workflow
lo usa, mantenere gli aggiornamenti automatizzati e non includere
credenziali del registry nel progetto. Configurazioni di feed e token
devono stare in ambienti o secret store protetti. La pipeline dovrebbe
ripristinare, compilare e testare da una copia pulita del repository.

\subsection{5.2.7 Un primo servizio e le verifiche di
ambiente}\label{un-primo-servizio-e-le-verifiche-di-ambiente}

Creare un progetto è solo la prima verifica. Prima di migrare logica
applicativa, accertarsi che il team possa:

\begin{enumerate}
\def\labelenumi{\arabic{enumi}.}
\tightlist
\item
  installare o selezionare la versione SDK prevista;
\item
  ripristinare i pacchetti dietro le restrizioni di rete del progetto;
\item
  compilare e testare dalla CLI e in CI;
\item
  avviare l'applicazione con la configurazione locale;
\item
  produrre un artefatto pubblicabile e avviarlo in un ambiente simile
  alla produzione.
\end{enumerate}

Registrare i comandi funzionanti e le versioni nel README del progetto.
Un ``funziona sul mio IDE'' non è ancora una toolchain condivisa.

\subsection{5.2.8 La mappa deve essere verificata sul caso
reale}\label{la-mappa-deve-essere-verificata-sul-caso-reale}

Per ogni concetto del sistema sorgente, compilare una scheda con tre
colonne: responsabilità attuale, candidato .NET, differenza da provare.
Per esempio: un hook di richiesta in un framework Python può essere
middleware, filtro o endpoint filter; la scelta dipende dal fatto che
debba avvolgere l'intera richiesta, agire solo su una route o validare
parametri. Questa verifica evita migrazioni basate sulla sola
somiglianza dei nomi.

\section{5.3 Pianificare una migrazione
controllata}\label{pianificare-una-migrazione-controllata}

Scelta la direzione, bisogna decidere come raggiungerla. La strategia
influenza tempi, costi, possibilità di rollback, coordinamento tra team
e probabilità di mantenere il servizio disponibile. Non esiste una
tecnica universalmente corretta: la decisione dipende dai confini
funzionali, dalle dipendenze condivise, dalle capacità di test e dai
vincoli del business.

\subsection{5.3.1 Tre strategie, non due}\label{tre-strategie-non-due}

\textbf{Riscrittura completa (big bang).} Il sistema viene ricostruito
in un ramo o ambiente separato e sostituito in un'unica data. È semplice
da spiegare, ma raramente semplice da eseguire: per molto tempo due
implementazioni divergono e il team scopre tardi differenze nei dati e
nei comportamenti. Ha senso solo con superficie contenuta, confini
chiari, test affidabili, breve durata e un cutover reversibile. Un lungo
congelamento funzionale può ridurre la divergenza, ma comporta costi di
business.

\textbf{Migrazione per componenti.} Si selezionano servizi o moduli
autonomi e li si migra in sequenza. Riduce il perimetro di ogni
rilascio, a condizione che i confini siano reali. Un modulo solo
nominalmente separato ma accoppiato a tabelle condivise, sessioni o
chiamate interne non è necessariamente una buona unità di migrazione.

\textbf{Sostituzione incrementale (\emph{strangler}).} Un router,
gateway o livello di ingresso dirige alcune capacità verso il sistema
nuovo e mantiene le restanti sul sistema originale. Le funzionalità
vengono spostate progressivamente; alla fine si rimuove il vecchio
sistema. Consente feedback reale e rollback selettivo, ma richiede di
gestire coesistenza, identità, osservabilità e contratti tra i due lati.
Non è sinonimo di riscrittura in microservizi.

{\def\LTcaptype{none} % do not increment counter
\begin{longtable}[]{@{}
  >{\raggedright\arraybackslash}p{(\linewidth - 6\tabcolsep) * \real{0.2500}}
  >{\raggedright\arraybackslash}p{(\linewidth - 6\tabcolsep) * \real{0.2500}}
  >{\raggedright\arraybackslash}p{(\linewidth - 6\tabcolsep) * \real{0.2500}}
  >{\raggedright\arraybackslash}p{(\linewidth - 6\tabcolsep) * \real{0.2500}}@{}}
\toprule\noalign{}
\begin{minipage}[b]{\linewidth}\raggedright
Strategia
\end{minipage} & \begin{minipage}[b]{\linewidth}\raggedright
Rischio concentrato
\end{minipage} & \begin{minipage}[b]{\linewidth}\raggedright
Vantaggio
\end{minipage} & \begin{minipage}[b]{\linewidth}\raggedright
Condizione che la rende praticabile
\end{minipage} \\
\midrule\noalign{}
\endhead
\bottomrule\noalign{}
\endlastfoot
Big bang & Cutover e regressioni aggregate & Confine di transizione
semplice & Sistema piccolo e test completo \\
Per componenti & Integrazione tra moduli & Progressi e rollback per
unità & Confini e ownership chiari \\
Strangler & Coesistenza e routing & Migrazione graduale in produzione &
Ingresso instradabile e contratti stabili \\
\end{longtable}
}

Una strategia ibrida è frequente: si crea un nuovo edge applicativo, si
migrano prima funzionalità isolate, si sostituiscono poi i flussi con
dipendenze più forti e si pianifica una finestra di cutover per elementi
non separabili.

\subsection{5.3.2 Scegliere la prima
capacità}\label{scegliere-la-prima-capacituxe0}

La prima funzionalità migrata non deve essere né un semplice endpoint
senza valore, né il flusso più critico dell'azienda. È preferibile una
capacità rappresentativa ma reversibile: possiede un contratto
comprensibile, ha dipendenze sotto controllo e permette di esercitare
build, deploy, logging, test e routing.

Valutare ogni candidato considerando valore, rischio, dipendenze, volume
di traffico e facilità di verifica. Un indice sintetico può aiutare a
confrontare opzioni, ma non deve sostituire la discussione: un punteggio
numerico non elimina l'incertezza. Il team dovrebbe poter spiegare
perché una capacità è pronta e quali evidenze mancano.

Prima di spostare il traffico occorre sapere come il router identifica
la funzionalità: prefisso di route, host, tenant, feature flag o regola
di gateway. Se la decisione dipende da un flag, deve essere possibile
disabilitarlo rapidamente senza ricompilare o perdere tracciabilità
delle modifiche. Un flag che resta indefinitamente nel codice non è un
piano di migrazione.

\subsection{5.3.3 Compatibilità durante la
coesistenza}\label{compatibilituxe0-durante-la-coesistenza}

Due sistemi attivi contemporaneamente possono leggere e scrivere lo
stesso dato. Le letture sono spesso più facili da spostare; le scritture
richiedono un proprietario chiaro. Scrivere su due database indipendenti
in sequenza non crea una transazione atomica: una scrittura può riuscire
e l'altra fallire. La doppia scrittura non va introdotta come soluzione
rapida senza un protocollo di riconciliazione, idempotenza e gestione
dei conflitti.

Le alternative dipendono dal dominio:

\begin{itemize}
\tightlist
\item
  un solo sistema possiede temporaneamente la scrittura; l'altro legge
  dati replicati;
\item
  si trasferisce la proprietà per entità o gruppo, con un registro di
  ownership;
\item
  si usa un outbox transazionale per pubblicare cambiamenti dopo il
  commit;
\item
  si effettua una migrazione con snapshot iniziale e sincronizzazione
  incrementale.
\end{itemize}

Ogni opzione ha costi e failure mode. L'outbox evita di perdere
l'intenzione di pubblicare un evento nella stessa transazione locale, ma
non rende atomici database diversi e consumer. La consegna può essere
ripetuta; i consumer devono essere idempotenti. In tutte le opzioni
vanno definiti conflitti, retry, deduplicazione e procedura di
riconciliazione.

Anche i contratti devono evolvere con compatibilità. Aggiungere un campo
opzionale è spesso più sicuro che rinominare o cambiare il tipo di uno
esistente; tuttavia il client potrebbe rifiutare campi sconosciuti o
interpretare diversamente un default. Versionamento e compatibilità
vanno provati contro i client reali o le loro specifiche.

\subsection{5.3.4 Fasi con condizioni di
uscita}\label{fasi-con-condizioni-di-uscita}

Una sequenza utile è:

\begin{enumerate}
\def\labelenumi{\arabic{enumi}.}
\tightlist
\item
  \textbf{Discovery:} inventario, baseline, dipendenze e rischi
  approvati.
\item
  \textbf{Foundation:} SDK/CI, hosting, log, configurazione, sicurezza e
  deploy vuoto verificati.
\item
  \textbf{Caratterizzazione:} contratti ed esempi rappresentativi
  fissati nei test.
\item
  \textbf{Vertical slice:} un flusso completo attraversa HTTP, dominio,
  persistenza, test e produzione non critica.
\item
  \textbf{Migrazione per capacità:} ogni unità ha criteri
  d'accettazione, osservazione e rollback.
\item
  \textbf{Cutover:} traffico spostato con una finestra e responsabili
  espliciti.
\item
  \textbf{Stabilizzazione e rimozione:} dati riconciliati, dipendenze
  eliminate, vecchio sistema dismesso e costi ridotti.
\end{enumerate}

Le fasi possono sovrapporsi, ma le condizioni d'uscita devono essere
esplicite. Per esempio, la foundation non è completata se esiste
soltanto un Dockerfile: bisogna dimostrare che l'immagine gira
nell'ambiente previsto, riceve configurazione esterna, espone health
check e termina correttamente.

\subsection{5.3.5 Cutover e rollback}\label{cutover-e-rollback}

Il cutover va trattato come un'operazione controllata, non come una
modifica improvvisata al load balancer. Il piano include orario,
responsabile della decisione, metriche da osservare, soglie d'arresto,
canali di comunicazione, ownership dei dati e comandi di ritorno.

Un rollback tecnico non annulla automaticamente le scritture già fatte
dal nuovo sistema. Se lo schema o il formato dati non sono leggibili dal
vecchio sistema, tornare indietro può richiedere restore, replay o
riconciliazione. Applicare migrazioni database compatibili in più passi:
espandere lo schema, distribuire codice che supporta entrambi i formati,
migrare i dati, poi rimuovere gli elementi obsoleti solo quando il
rollback non li richiede più.

Provare il rollback prima del cutover. Registrare durata, autorizzazioni
necessarie, effetti sui dati e segnali che confermano che il traffico è
tornato sul vecchio percorso. Un piano non provato è un'ipotesi.

\subsection{5.3.6 Organizzazione e
dipendenze}\label{organizzazione-e-dipendenze}

La migrazione richiede decisioni condivise tra sviluppo, piattaforma,
sicurezza, database, supporto e product owner. Stabilire chi approva un
contratto, chi controlla le metriche e chi può fermare il rollout.
Esplicitare le dipendenze tra attività, evitando che più team cambino lo
stesso schema o endpoint senza coordinamento.

Il backlog deve distinguere lavoro di migrazione, correzioni necessarie
per la parità e nuove funzionalità. Questa separazione rende visibile il
costo e impedisce che una funzionalità aggiunta a metà diventi la
ragione per rinviare la dismissione. Ogni milestone dovrebbe produrre un
incremento funzionante e una decisione informata sul passo successivo.

\section{5.4 Architettura e organizzazione del
codice}\label{architettura-e-organizzazione-del-codice}

L'architettura di destinazione non dovrebbe essere una copia dei package
del framework sorgente né una collezione di pattern scelti per moda. Il
suo compito è rendere espliciti i confini: dove termina HTTP, dove
risiede la logica di dominio, chi possiede i dati e come si integrano le
dipendenze esterne.

\subsection{5.4.1 Dal framework sorgente alle responsabilità ASP.NET
Core}\label{dal-framework-sorgente-alle-responsabilituxe0-asp.net-core}

Laravel, Django, Express/NestJS e Spring Boot offrono combinazioni
diverse di routing, middleware, controller, container, configurazione e
persistenza. Durante la migrazione, decomporre il comportamento:

{\def\LTcaptype{none} % do not increment counter
\begin{longtable}[]{@{}
  >{\raggedright\arraybackslash}p{(\linewidth - 2\tabcolsep) * \real{0.5000}}
  >{\raggedright\arraybackslash}p{(\linewidth - 2\tabcolsep) * \real{0.5000}}@{}}
\toprule\noalign{}
\begin{minipage}[b]{\linewidth}\raggedright
Responsabilità sorgente
\end{minipage} & \begin{minipage}[b]{\linewidth}\raggedright
Possibile destinazione ASP.NET Core
\end{minipage} \\
\midrule\noalign{}
\endhead
\bottomrule\noalign{}
\endlastfoot
Route/controller & Minimal API o controller MVC \\
Hook globale per richiesta & Middleware \\
Hook limitato a endpoint & Endpoint filter o filtro MVC \\
Container dei servizi & \texttt{IServiceCollection} e provider DI \\
Configurazione ambientale & Configuration providers e Options \\
Worker/job & \texttt{BackgroundService}, worker o scheduler esterno \\
Validazione input & Validazione esplicita, filtri o libreria adottata \\
Template/server rendering & Razor Pages, MVC o frontend separato \\
\end{longtable}
}

La tabella indica candidati, non una traduzione automatica. Un
middleware ASP.NET Core avvolge una richiesta e può decidere se
proseguire; un filtro MVC opera nel ciclo MVC; un endpoint filter si
applica ai delegate degli endpoint. Scegliere il livello in base alla
responsabilità evita middleware globali che contengono eccezioni
specifiche per una sola route.

\subsection{5.4.2 Minimal API o
controller?}\label{minimal-api-o-controller}

Le Minimal API sono un modo conciso per definire endpoint e si adattano
bene a servizi piccoli, API focalizzate e team che vogliono comporre
route per feature. Non sono un framework ``senza architettura'':
dipendenze, validazione, errori, mapping e test devono comunque essere
organizzati.

I controller MVC offrono convenzioni e struttura familiare per API
articolate, applicazioni con molti endpoint simili, filtri MVC e team
che beneficiano di classi e attributi. Non obbligano a usare tutte le
funzionalità MVC. La scelta non va basata sul numero di righe iniziali:
valutare complessità delle regole trasversali, standardizzazione,
familiarità del team e necessità di organizzare i contratti.

Esempio Minimal API con gruppo:

\begin{Shaded}
\begin{Highlighting}[]
\DataTypeTok{var}\NormalTok{ builder }\OperatorTok{=}\NormalTok{ WebApplication}\OperatorTok{.}\FunctionTok{CreateBuilder}\OperatorTok{(}\NormalTok{args}\OperatorTok{);}
\NormalTok{builder}\OperatorTok{.}\FunctionTok{Services}\OperatorTok{.}\FunctionTok{AddScoped}\OperatorTok{\textless{}}\NormalTok{IProductService}\OperatorTok{,}\NormalTok{ ProductService}\OperatorTok{\textgreater{}();}

\DataTypeTok{var}\NormalTok{ app }\OperatorTok{=}\NormalTok{ builder}\OperatorTok{.}\FunctionTok{Build}\OperatorTok{();}
\DataTypeTok{var}\NormalTok{ products }\OperatorTok{=}\NormalTok{ app}\OperatorTok{.}\FunctionTok{MapGroup}\OperatorTok{(}\StringTok{"/api/products"}\OperatorTok{);}

\NormalTok{products}\OperatorTok{.}\FunctionTok{MapGet}\OperatorTok{(}\StringTok{"/\{id:int\}"}\OperatorTok{,} \FunctionTok{async} \OperatorTok{(}
    \DataTypeTok{int}\NormalTok{ id}\OperatorTok{,}\NormalTok{ IProductService service}\OperatorTok{,}\NormalTok{ CancellationToken ct}\OperatorTok{)} \OperatorTok{=\textgreater{}}
\OperatorTok{\{}
    \DataTypeTok{var}\NormalTok{ product }\OperatorTok{=}\NormalTok{ await service}\OperatorTok{.}\FunctionTok{FindAsync}\OperatorTok{(}\NormalTok{id}\OperatorTok{,}\NormalTok{ ct}\OperatorTok{);}
    \KeywordTok{return}\NormalTok{ product }\KeywordTok{is} \KeywordTok{null} \OperatorTok{?}\NormalTok{ Results}\OperatorTok{.}\FunctionTok{NotFound}\OperatorTok{()} \OperatorTok{:}\NormalTok{ Results}\OperatorTok{.}\FunctionTok{Ok}\OperatorTok{(}\NormalTok{product}\OperatorTok{);}
\OperatorTok{\});}

\NormalTok{app}\OperatorTok{.}\FunctionTok{Run}\OperatorTok{();}
\end{Highlighting}
\end{Shaded}

Il delegate riceve parametri che ASP.NET Core associa a route, servizi o
body secondo regole di binding. Per mantenere leggibile il punto di
composizione, spostare la definizione in un metodo
\texttt{MapProductEndpoints} o in un modulo quando le route crescono.
Non trasformare subito ogni gruppo in un'astrazione generica: estrarre
quando si chiarisce una responsabilità riutilizzabile.

\subsection{5.4.3 Confini applicativi e
dipendenze}\label{confini-applicativi-e-dipendenze}

Una struttura minima può separare:

\begin{itemize}
\tightlist
\item
  \textbf{API:} protocollo HTTP, autenticazione, mapping
  request/response;
\item
  \textbf{Application:} casi d'uso e coordinamento;
\item
  \textbf{Domain:} regole e invarianti di business;
\item
  \textbf{Infrastructure:} database e integrazioni esterne.
\end{itemize}

Questi nomi non sono obbligatori. In un servizio piccolo, più progetti
possono aggiungere overhead senza proteggere alcun confine. In un
sistema con dominio complesso e più adapter, la separazione riduce
l'accoppiamento. La direzione desiderata è che le regole del dominio non
dipendano dal server web o dal provider database; i dettagli
infrastrutturali implementano contratti necessari ai casi d'uso.

Un layout possibile per feature:

\begin{Shaded}
\begin{Highlighting}[]
\NormalTok{src/}
\NormalTok{  Catalogo.Api/}
\NormalTok{    Program.cs}
\NormalTok{    Features/}
\NormalTok{      Products/}
\NormalTok{        ProductEndpoints.cs}
\NormalTok{        CreateProductRequest.cs}
\NormalTok{  Catalogo.Application/}
\NormalTok{  Catalogo.Domain/}
\NormalTok{  Catalogo.Infrastructure/}
\NormalTok{tests/}
\NormalTok{  Catalogo.UnitTests/}
\NormalTok{  Catalogo.IntegrationTests/}
\end{Highlighting}
\end{Shaded}

Questo è uno schema, non un requisito. Una cartella
\texttt{Features/Products} può contenere endpoint e request specifiche,
mentre servizi applicativi e dominio vivono in progetti distinti.
Evitare cartelle generiche \texttt{Helpers}, \texttt{Common} o
\texttt{Utils} come destinazione per qualsiasi codice non collocato:
diventano rapidamente dipendenze senza ownership.

\subsection{5.4.4 Dominio, DTO e mapping}\label{dominio-dto-e-mapping}

Non esporre direttamente l'entità database come contratto pubblico per
comodità. Una modifica interna allo schema potrebbe cambiare
involontariamente l'API; inoltre l'entità può contenere campi che il
client non deve vedere. Usare DTO di richiesta e risposta quando esiste
un confine reale, definendo mapping espliciti per campi, null e valori
predefiniti.

\begin{Shaded}
\begin{Highlighting}[]
\KeywordTok{public} \KeywordTok{sealed}\NormalTok{ record }\FunctionTok{ProductResponse}\OperatorTok{(}
    \DataTypeTok{int}\NormalTok{ Id}\OperatorTok{,} \DataTypeTok{string}\NormalTok{ Name}\OperatorTok{,} \DataTypeTok{decimal}\NormalTok{ Price}\OperatorTok{,} \DataTypeTok{string}\NormalTok{ Currency}\OperatorTok{);}
\end{Highlighting}
\end{Shaded}

La logica di conversione deve essere testata. Un mapping automatico può
essere utile, ma non deve nascondere conversioni lossy, date senza
timezone o proprietà omesse. Per ogni contratto documentare proprietà
obbligatorie, formato, unità, precisione e compatibilità.

\subsection{5.4.5 Hosting e Kestrel}\label{hosting-e-kestrel}

\texttt{WebApplication.CreateBuilder(args)} compone configuration,
logging, servizi e ambiente. L'app risultante può girare in foreground,
in container o come servizio gestito. Kestrel gestisce HTTP; in
produzione spesso c'è un reverse proxy o un ingress davanti. Questo
proxy può terminare TLS, imporre limiti, fare routing e aggiungere
header forwarded. L'app deve fidarsi solo dei proxy configurati, non di
qualunque header fornito dal client.

Il binding delle porte e degli URL deve provenire dalla configurazione
dell'ambiente. Non codificare \texttt{localhost} o porte di sviluppo
nell'immagine. Verificare che il processo ascolti sull'interfaccia
prevista, che i controlli health siano raggiungibili dal sistema di
orchestrazione e che lo shutdown consenta di terminare richieste in
corso.

\subsection{5.4.6 Processi web, worker e
job}\label{processi-web-worker-e-job}

Un'applicazione sorgente può eseguire job nello stesso processo HTTP. In
.NET un \texttt{BackgroundService} è appropriato per lavoro persistente
controllato dall'host, ma non è una coda durevole: se il processo si
arresta, gli elementi in memoria possono andare persi. Per job che
richiedono retry, scheduling distribuito o garanzie di consegna, usare
una coda o un sistema di scheduling adeguato.

I worker devono gestire cancellazione e shutdown, limitare concorrenza,
registrare gli errori e rendere idempotenti le operazioni ripetibili.
Non avviare thread detached da un endpoint e non ignorare eccezioni in
task non osservati. Il lifecycle del lavoro deve essere parte della
strategia di deployment e scaling.

\subsection{5.4.7 Strutturare per comprendere, non per
imitare}\label{strutturare-per-comprendere-non-per-imitare}

Portare il modello mentale del vecchio framework è utile come punto di
partenza, ma il target deve rispettare le sue convenzioni. Per ogni
astrazione chiedersi: quale responsabilità protegge? Quale test la rende
verificabile? Quale cambiamento futuro diventa più semplice? Se non c'è
una risposta, l'astrazione probabilmente anticipa un problema non ancora
dimostrato.

Una migrazione ben strutturata è quella che permette di trovare
rapidamente un endpoint, capire chi possiede una regola, sostituire una
dipendenza esterna e testare il comportamento senza avviare l'intero
sistema. Non è necessariamente quella con più progetti o livelli.

\section{5.5 Richieste HTTP e pipeline}\label{richieste-http-e-pipeline}

Un'API migra con successo solo se i suoi client osservano il contratto
previsto. Cambiare framework spesso cambia dettagli che non erano
scritti: codici di stato, forma degli errori, parsing di date,
comportamento sui campi null, header, encoding e ordine delle
operazioni. Questa sezione tratta la pipeline HTTP come un contratto da
preservare e testare.

\subsection{5.5.1 Ricostruire il contratto prima di spostare le
route}\label{ricostruire-il-contratto-prima-di-spostare-le-route}

Per ciascun endpoint registrare metodo e percorso, autenticazione,
parametri, schema request/response, codici di stato, header,
paginazione, ordinamento, idempotenza e limiti. Salvare esempi
realistici, inclusi input invalidi e casi limite. OpenAPI aiuta a
descrivere il contratto ma non sempre riflette il comportamento
effettivo; confrontare specifica, implementazione e traffico
anonimizzato.

Attenzione alle differenze di parsing:

\begin{itemize}
\tightlist
\item
  distinzione maiuscole/minuscole nei nomi JSON e nei valori;
\item
  date locali, offset e precisione dei timestamp;
\item
  separatori decimali e formati numerici;
\item
  \texttt{null}, proprietà assente, stringa vuota e valore predefinito;
\item
  encoding, Unicode normalizzato e limiti sulla dimensione del body;
\item
  parametri duplicati, query string e codifica degli spazi.
\end{itemize}

Una modifica ``più corretta'' può essere breaking. Se si decide di
correggere un comportamento, documentare la variazione e verificare la
compatibilità dei consumer.

\subsection{5.5.2 Middleware e ordine della
pipeline}\label{middleware-e-ordine-della-pipeline}

Il middleware è una sequenza. Un componente può modificare la richiesta
prima del successivo, osservare la risposta dopo \texttt{await\ next()},
oppure terminare la pipeline. L'ordine influenza error handling,
routing, autenticazione e autorizzazione.

\begin{Shaded}
\begin{Highlighting}[]
\DataTypeTok{var}\NormalTok{ builder }\OperatorTok{=}\NormalTok{ WebApplication}\OperatorTok{.}\FunctionTok{CreateBuilder}\OperatorTok{(}\NormalTok{args}\OperatorTok{);}
\NormalTok{builder}\OperatorTok{.}\FunctionTok{Services}\OperatorTok{.}\FunctionTok{AddProblemDetails}\OperatorTok{();}
\NormalTok{builder}\OperatorTok{.}\FunctionTok{Services}\OperatorTok{.}\FunctionTok{AddControllers}\OperatorTok{();}
\NormalTok{builder}\OperatorTok{.}\FunctionTok{Services}\OperatorTok{.}\FunctionTok{AddAuthentication}\OperatorTok{();}
\NormalTok{builder}\OperatorTok{.}\FunctionTok{Services}\OperatorTok{.}\FunctionTok{AddAuthorization}\OperatorTok{();}

\DataTypeTok{var}\NormalTok{ app }\OperatorTok{=}\NormalTok{ builder}\OperatorTok{.}\FunctionTok{Build}\OperatorTok{();}

\NormalTok{app}\OperatorTok{.}\FunctionTok{UseExceptionHandler}\OperatorTok{();}
\NormalTok{app}\OperatorTok{.}\FunctionTok{UseHttpsRedirection}\OperatorTok{();}
\NormalTok{app}\OperatorTok{.}\FunctionTok{UseRouting}\OperatorTok{();}
\NormalTok{app}\OperatorTok{.}\FunctionTok{UseAuthentication}\OperatorTok{();}
\NormalTok{app}\OperatorTok{.}\FunctionTok{UseAuthorization}\OperatorTok{();}
\NormalTok{app}\OperatorTok{.}\FunctionTok{MapControllers}\OperatorTok{();}
\NormalTok{app}\OperatorTok{.}\FunctionTok{Run}\OperatorTok{();}
\end{Highlighting}
\end{Shaded}

Per rendere operativo il gestore standard, registrare prima
\texttt{AddProblemDetails}. In assenza di un handler personalizzato, il
middleware restituisce una risposta Problem Details quando il servizio è
disponibile. L'ordine reale dipende da servizi e hosting: seguire la
documentazione della versione ASP.NET Core scelta. In generale il
gestore eccezioni deve avvolgere i componenti che si vogliono gestire;
autenticazione deve precedere autorizzazione; il routing deve aver
stabilito l'endpoint prima che componenti dipendenti dall'endpoint lo
usino. Non copiare meccanicamente un ordine trovato in un progetto
diverso.

Un middleware per correlare le richieste deve propagare un
identificatore sicuro e aggiungere contesto ai log, ma non deve fidarsi
ciecamente di un valore arbitrario passato dall'esterno. Anche un header
di correlazione può essere usato per iniettare caratteri o dati
ingannevoli nei log; normalizzare e limitare la lunghezza.

\subsection{5.5.3 Routing e binding}\label{routing-e-binding}

Minimal API e controller offrono convenzioni diverse per dichiarare
route e binding. Una route esplicita rende visibile il contratto:

\begin{Shaded}
\begin{Highlighting}[]
\NormalTok{app}\OperatorTok{.}\FunctionTok{MapGet}\OperatorTok{(}\StringTok{"/api/products/\{id:int\}"}\OperatorTok{,} \FunctionTok{async} \OperatorTok{(}
    \DataTypeTok{int}\NormalTok{ id}\OperatorTok{,}\NormalTok{ IProductService service}\OperatorTok{,}\NormalTok{ CancellationToken ct}\OperatorTok{)} \OperatorTok{=\textgreater{}}
\OperatorTok{\{}
    \DataTypeTok{var}\NormalTok{ item }\OperatorTok{=}\NormalTok{ await service}\OperatorTok{.}\FunctionTok{FindAsync}\OperatorTok{(}\NormalTok{id}\OperatorTok{,}\NormalTok{ ct}\OperatorTok{);}
    \KeywordTok{return}\NormalTok{ item }\KeywordTok{is} \KeywordTok{null}
        \OperatorTok{?}\NormalTok{ Results}\OperatorTok{.}\FunctionTok{NotFound}\OperatorTok{()}
        \OperatorTok{:}\NormalTok{ Results}\OperatorTok{.}\FunctionTok{Ok}\OperatorTok{(}\NormalTok{item}\OperatorTok{);}
\OperatorTok{\});}
\end{Highlighting}
\end{Shaded}

La constraint \texttt{:int} restringe il matching; non sostituisce una
risposta coerente per input non valido. Per request complesse definire
un tipo DTO anziché legare un'entità interna. Evitare parametri
implicitamente ambigui e testare l'overload di route, slash finali,
versioni e route non trovate.

La validazione deve avvenire al confine e includere vincoli di dominio
pertinenti. Un controllo sintattico (campo obbligatorio, lunghezza
massima) è diverso da una regola di business (prodotto attivo, prezzo
modificabile). Restituire errori leggibili e stabili, senza stack trace
o dettagli di database.

\subsection{5.5.4 Errori prevedibili}\label{errori-prevedibili}

Definire una tassonomia di errori: input non valido, risorsa
inesistente, conflitto, non autenticato, non autorizzato, dipendenza non
disponibile e errore inatteso. Mappare ogni categoria a status code e
corpo consistente. \texttt{ProblemDetails} è un formato standard utile
per risposte HTTP, ma i dettagli sensibili devono essere esclusi in
produzione.

Le eccezioni non dovrebbero essere usate come controllo di flusso
ordinario. La logica applicativa può restituire un risultato esplicito
per conflitti attesi; il middleware centrale gestisce gli errori
inattesi, registra l'eccezione e restituisce un identificatore di
correlazione. Non restituire HTTP 200 con
\texttt{\{\ "success":\ false\ \}} se il contratto prevede status
semantici; non trasformare ogni errore in 500 senza classificazione.

Un \texttt{catch} locale è appropriato se il componente può recuperare,
aggiungere contesto utile o tradurre una specifica eccezione in un
risultato noto. Altrimenti lasciar propagare al gestore centrale.
Evitare log duplicati della stessa eccezione a ogni livello.

\subsection{5.5.5 Filtri, middleware ed endpoint
filter}\label{filtri-middleware-ed-endpoint-filter}

La scelta del punto d'estensione dipende dall'ambito:

\begin{itemize}
\tightlist
\item
  \textbf{middleware:} cross-cutting a livello di pipeline, per log,
  header o gestione globale;
\item
  \textbf{filtro MVC:} comportamento nel lifecycle dei controller, per
  endpoint MVC;
\item
  \textbf{endpoint filter:} logica riutilizzabile per gruppi di Minimal
  API, inclusi controlli su argomenti;
\item
  \textbf{codice applicativo:} regole specifiche del caso d'uso che
  devono essere testabili senza HTTP.
\end{itemize}

Non collocare regole di dominio in middleware solo perché è il primo
hook conosciuto. Un middleware non dovrebbe conoscere il concetto di
``ordine annullato'' o ``cliente premium''.

\subsection{5.5.6 Timeout, cancellazione e
resilienza}\label{timeout-cancellazione-e-resilienza}

Ogni chiamata remota può fallire o impiegare troppo tempo. Propagare
\texttt{CancellationToken} dalle API alle query e ai client HTTP;
applicare timeout espliciti e limiti per componente. Un timeout del
client non garantisce che l'operazione remota sia stata annullata:
l'esito può essere ambiguo, quindi le operazioni di scrittura devono
considerare idempotency key o meccanismi di verifica.

Retry automatici sono adatti solo a errori transitori e operazioni
sicure da ripetere. Non fare retry indiscriminati su ogni eccezione o su
operazioni non idempotenti. Usare backoff con jitter, budget
complessivo, circuit breaker se utile e telemetria che mostri i
tentativi. Evitare retry annidati in più livelli: moltiplicano la
pressione sul servizio già degradato.

\subsection{5.5.7 CORS, rate limiting e
documentazione}\label{cors-rate-limiting-e-documentazione}

CORS è una policy applicata dai browser, non un'autenticazione.
Consentire origini, metodi e header minimi; non combinare wildcard con
credenziali. Il rate limiting protegge risorse e client equamente, ma
una policy in memoria per istanza non produce limiti globali in un
cluster. Scegliere chiave e storage coerenti con il deployment,
proteggendo le chiavi da spoofing.

OpenAPI descrive endpoint e schemi e può generare client, ma la
specifica deve essere parte della revisione. Mantenere esempi, errori,
autenticazione e versionamento aggiornati. Diff automatizzati tra
specifiche aiutano a individuare breaking change prima del rilascio.

\subsection{5.5.8 Checklist della parità
HTTP}\label{checklist-della-parituxe0-http}

Prima di spostare il traffico verificare: status code; schema e casing
JSON; header; autenticazione; codici di errore; date e timezone;
paginazione; ordinamento; CORS; timeout; payload grandi; input
malformati; richieste duplicate; cancellazione e versionamento. Eseguire
i test contro vecchio e nuovo sistema con gli stessi casi, confrontando
gli output normalizzati e documentando ogni differenza intenzionale.

\section{5.6 Configurazione, stato e
dipendenze}\label{configurazione-stato-e-dipendenze}

I framework di partenza spesso rendono facili da condividere servizi,
configurazione e stato di richiesta. In ASP.NET Core queste
responsabilità sono esplicite: il container DI crea le dipendenze, i
provider di configurazione aggregano valori e i servizi di cache o
sessione gestiscono lo stato. Esplicitarle è un vantaggio solo se il
lifecycle scelto è corretto.

\subsection{5.6.1 Dependency Injection
integrata}\label{dependency-injection-integrata}

La registrazione avviene prima di costruire l'applicazione:

\begin{Shaded}
\begin{Highlighting}[]
\DataTypeTok{var}\NormalTok{ builder }\OperatorTok{=}\NormalTok{ WebApplication}\OperatorTok{.}\FunctionTok{CreateBuilder}\OperatorTok{(}\NormalTok{args}\OperatorTok{);}
\NormalTok{builder}\OperatorTok{.}\FunctionTok{Services}\OperatorTok{.}\FunctionTok{AddScoped}\OperatorTok{\textless{}}\NormalTok{IOrderService}\OperatorTok{,}\NormalTok{ OrderService}\OperatorTok{\textgreater{}();}
\NormalTok{builder}\OperatorTok{.}\FunctionTok{Services}\OperatorTok{.}\FunctionTok{AddSingleton}\OperatorTok{\textless{}}\NormalTok{IClock}\OperatorTok{,}\NormalTok{ SystemClock}\OperatorTok{\textgreater{}();}
\NormalTok{builder}\OperatorTok{.}\FunctionTok{Services}\OperatorTok{.}\FunctionTok{AddHttpClient}\OperatorTok{\textless{}}\NormalTok{PaymentClient}\OperatorTok{\textgreater{}();}
\DataTypeTok{var}\NormalTok{ app }\OperatorTok{=}\NormalTok{ builder}\OperatorTok{.}\FunctionTok{Build}\OperatorTok{();}
\end{Highlighting}
\end{Shaded}

Il container integrato risolve i servizi e li libera secondo il
lifecycle. I tre lifetime principali sono:

\begin{itemize}
\tightlist
\item
  \textbf{Transient:} nuova istanza a ogni risoluzione. Adatto a oggetti
  leggeri e senza stato condiviso.
\item
  \textbf{Scoped:} una istanza per scope; nelle API web di norma uno
  scope per richiesta. \texttt{DbContext} è tipicamente scoped.
\item
  \textbf{Singleton:} una istanza per l'intera durata del processo. Deve
  essere thread-safe e non conservare dati specifici della richiesta.
\end{itemize}

Il lifetime non definisce soltanto l'efficienza. Un singleton che
dipende da uno scoped può catturare un oggetto oltre la sua durata
(\emph{captive dependency}). Un servizio scoped risolto dal singleton
può accedere a un \texttt{DbContext} condiviso in modo non sicuro; due
richieste concorrenti potrebbero interferire. L'abilitazione delle
verifiche di scope in sviluppo aiuta a rilevare configurazioni errate,
ma i test devono coprire il comportamento reale.

Iniettare attraverso il costruttore rende visibili le dipendenze:

\begin{Shaded}
\begin{Highlighting}[]
\KeywordTok{public} \KeywordTok{sealed} \KeywordTok{class} \FunctionTok{OrderService}\OperatorTok{(}
\NormalTok{    IOrderRepository repository}\OperatorTok{,}
\NormalTok{    IClock clock}\OperatorTok{)} \OperatorTok{:}\NormalTok{ IOrderService}
\OperatorTok{\{}
    \KeywordTok{public}\NormalTok{ Task}\OperatorTok{\textless{}}\NormalTok{Order}\OperatorTok{?\textgreater{}} \FunctionTok{FindAsync}\OperatorTok{(}
        \DataTypeTok{int}\NormalTok{ id}\OperatorTok{,}\NormalTok{ CancellationToken ct}\OperatorTok{)} \OperatorTok{=\textgreater{}}
\NormalTok{        repository}\OperatorTok{.}\FunctionTok{FindAsync}\OperatorTok{(}\NormalTok{id}\OperatorTok{,}\NormalTok{ ct}\OperatorTok{);}
\OperatorTok{\}}
\end{Highlighting}
\end{Shaded}

Una classe con moltissimi parametri può indicare che coordina troppe
responsabilità. Non risolvere il problema iniettando il provider
generico \texttt{IServiceProvider} e facendo service locator: le
dipendenze diventano invisibili, l'errore emerge a runtime e i test
richiedono più setup.

\subsection{5.6.2 Configuration providers e
priorità}\label{configuration-providers-e-priorituxe0}

La configurazione standard combina valori da più fonti, tipicamente file
JSON, variabili d'ambiente, user secrets in sviluppo e argomenti CLI. I
provider successivi possono sovrascrivere valori precedenti; conoscere
l'ordine è essenziale per capire quale valore è effettivamente attivo.

\begin{Shaded}
\begin{Highlighting}[]
\FunctionTok{\{}
  \DataTypeTok{"Catalog"}\FunctionTok{:} \FunctionTok{\{}
    \DataTypeTok{"PageSize"}\FunctionTok{:} \DecValTok{50}\FunctionTok{,}
    \DataTypeTok{"EnableBetaSearch"}\FunctionTok{:} \KeywordTok{false}
  \FunctionTok{\}}
\FunctionTok{\}}
\end{Highlighting}
\end{Shaded}

Le variabili d'ambiente usano convenzioni come
\texttt{Catalog\_\_PageSize=100} per rappresentare la gerarchia.
\texttt{appsettings.json} contiene configurazione non segreta e valori
predefiniti; non vi si inseriscono password, token o chiavi di
produzione. La configurazione non è una stringa da leggere liberamente
in ogni metodo: raggruppare opzioni correlate in tipi tipizzati,
validarne i valori all'avvio e fallire esplicitamente se manca un
requisito essenziale.

\begin{Shaded}
\begin{Highlighting}[]
\NormalTok{builder}\OperatorTok{.}\FunctionTok{Services}
    \OperatorTok{.}\FunctionTok{AddOptions}\OperatorTok{\textless{}}\NormalTok{CatalogOptions}\OperatorTok{\textgreater{}()}
    \OperatorTok{.}\FunctionTok{BindConfiguration}\OperatorTok{(}\StringTok{"Catalog"}\OperatorTok{)}
    \OperatorTok{.}\FunctionTok{Validate}\OperatorTok{(}\NormalTok{o }\OperatorTok{=\textgreater{}}\NormalTok{ o}\OperatorTok{.}\FunctionTok{PageSize} \KeywordTok{is} \OperatorTok{\textgreater{}} \DecValTok{0}\NormalTok{ and }\OperatorTok{\textless{}=} \DecValTok{200}\OperatorTok{,}
        \StringTok{"PageSize deve essere compreso tra 1 e 200."}\OperatorTok{)}
    \OperatorTok{.}\FunctionTok{ValidateOnStart}\OperatorTok{();}
\end{Highlighting}
\end{Shaded}

Un default silenzioso per una configurazione obbligatoria può produrre
un'applicazione che si avvia ma invia dati all'host sbagliato o
disabilita una protezione. Preferire un errore chiaro in fase di avvio e
loggare il nome della chiave mancante, non il suo valore segreto.

\subsection{5.6.3 Secrets e ambienti}\label{secrets-e-ambienti}

In sviluppo, User Secrets evita di commettere segreti nel repository, ma
non è un vault di produzione. In produzione, usare il secret manager
della piattaforma o un provider protetto, con rotazione, accesso minimo
e audit. Evitare segreti in immagini container, argomenti di processo e
log.

Separare configurazione per ambiente senza duplicare interi file con
valori divergenti. Il codice dovrebbe leggere la configurazione da
un'unica interfaccia; è l'ambiente a fornire il valore. La pipeline deve
poter distribuire lo stesso artefatto in ambienti diversi e applicare i
parametri esternamente.

Una chiave di produzione non deve essere copiata su una macchina locale
per ``far funzionare'' il test. Usare credenziali effimere, identità
gestite o servizi sostitutivi con permessi limitati.

\subsection{5.6.4 Stato: richiesta, sessione e
cache}\label{stato-richiesta-sessione-e-cache}

Lo stato di una richiesta HTTP vive nell'oggetto \texttt{HttpContext} e
non va conservato oltre la richiesta. Il server può servire richieste
concorrenti su più thread: campi statici o singleton mutabili non sono
un sostituto della sessione.

Distinguere:

\begin{itemize}
\tightlist
\item
  \textbf{stato di dominio persistente:} database o sistema autorevole;
\item
  \textbf{cache:} dato ricostruibile con scadenza, invalidazione e
  comportamento in caso di miss;
\item
  \textbf{sessione:} dati temporanei legati a un utente, con
  implicazioni di privacy e scalabilità;
\item
  \textbf{stato locale al processo:} utile per ottimizzazioni effimere,
  mai fonte di verità in un cluster.
\end{itemize}

Con più istanze, cache o sessioni in memoria divergono. Una richiesta
successiva può arrivare a un nodo diverso; sticky sessions riducono
alcuni sintomi ma non sostituiscono storage distribuito e resilienza.
Documentare durata, chiave, invalidazione, consistenza e comportamento
in caso di indisponibilità della cache.

Non mettere nella cache dati sensibili o per-utente usando chiavi non
isolate. Stabilire un limite di memoria e TTL; un cache stampede può
saturare database e servizi proprio quando una chiave scade. Le
strategie di refresh in background, lock distribuiti o
stale-while-revalidate vanno scelte in base alla tolleranza alla
staleness.

\subsection{5.6.5 Client HTTP e
integrazione}\label{client-http-e-integrazione}

Per chiamare servizi esterni usare \texttt{IHttpClientFactory} o client
tipizzati; evitare di creare e distruggere \texttt{HttpClient} a ogni
chiamata. Configurare base address, timeout, handler e credenziali in
modo centralizzato. Un client tipizzato rende visibili gli endpoint
esterni e facilita test, logging e policy di resilienza.

Non registrare token nei log. Configurare timeout sia per la richiesta
sia per operazioni specifiche; un timeout infinito può accumulare
richieste e socket. Distinguere errori DNS, connessione rifiutata,
timeout, status non riuscito e contenuto non valido. La risposta di un
sistema remoto è input non attendibile anche quando usa TLS.

\subsection{5.6.6 Background services e lavoro
durevole}\label{background-services-e-lavoro-durevole}

\texttt{BackgroundService} avvia un task con il lifecycle dell'host. È
adatto a polling o lavori periodici semplici, non è una coda
persistente. Se serve consegna affidabile, retry dopo crash, scheduling
distribuito o visibilità operativa, usare un broker o un job scheduler
con storage durevole.

Quando il worker usa servizi scoped, creare uno scope per l'unità di
lavoro:

\begin{Shaded}
\begin{Highlighting}[]
\KeywordTok{public} \KeywordTok{sealed} \KeywordTok{class} \FunctionTok{ImportWorker}\OperatorTok{(}
\NormalTok{    IServiceScopeFactory scopeFactory}\OperatorTok{,}
\NormalTok{    ILogger}\OperatorTok{\textless{}}\NormalTok{ImportWorker}\OperatorTok{\textgreater{}}\NormalTok{ logger}\OperatorTok{)} \OperatorTok{:}\NormalTok{ BackgroundService}
\OperatorTok{\{}
    \KeywordTok{protected} \KeywordTok{override}\NormalTok{ async Task }\FunctionTok{ExecuteAsync}\OperatorTok{(}
\NormalTok{        CancellationToken stoppingToken}\OperatorTok{)}
    \OperatorTok{\{}
        \KeywordTok{while} \OperatorTok{(!}\NormalTok{stoppingToken}\OperatorTok{.}\FunctionTok{IsCancellationRequested}\OperatorTok{)}
        \OperatorTok{\{}
\NormalTok{            await }\KeywordTok{using} \DataTypeTok{var}\NormalTok{ scope }\OperatorTok{=}\NormalTok{ scopeFactory}\OperatorTok{.}\FunctionTok{CreateAsyncScope}\OperatorTok{();}
            \DataTypeTok{var}\NormalTok{ importer }\OperatorTok{=}\NormalTok{ scope}\OperatorTok{.}\FunctionTok{ServiceProvider}
                \OperatorTok{.}\FunctionTok{GetRequiredService}\OperatorTok{\textless{}}\NormalTok{IImporter}\OperatorTok{\textgreater{}();}
\NormalTok{            await importer}\OperatorTok{.}\FunctionTok{ProcessBatchAsync}\OperatorTok{(}\NormalTok{stoppingToken}\OperatorTok{);}
\NormalTok{            await Task}\OperatorTok{.}\FunctionTok{Delay}\OperatorTok{(}\NormalTok{TimeSpan}\OperatorTok{.}\FunctionTok{FromSeconds}\OperatorTok{(}\DecValTok{10}\OperatorTok{),}\NormalTok{ stoppingToken}\OperatorTok{);}
        \OperatorTok{\}}
    \OperatorTok{\}}
\OperatorTok{\}}
\end{Highlighting}
\end{Shaded}

In codice reale gestire cancellazione e backoff: un'eccezione non
recuperata può terminare il worker o il processo secondo configurazione.
Il batch deve essere ripetibile, limitato e osservabile. Non avviare un
task fire-and-forget da una route, perché l'host non ne governa durata o
errori.

\subsection{5.6.7 Progettare le dipendenze come
confini}\label{progettare-le-dipendenze-come-confini}

Ogni dipendenza esterna dovrebbe essere descritta da un'interfaccia o da
un adapter quando questo aiuta a isolare un confine reale: provider di
pagamento, clock, file storage o API remota. Non serve creare
un'interfaccia per ogni classe. L'obiettivo è sostituire un'integrazione
durante test o migrazione senza nascondere le capacità utilizzate.

Annotare chi possiede la configurazione e il lifecycle. Nei test, usare
servizi in-memory soltanto se preservano la semantica necessaria; un
dizionario non simula transazioni o concorrenza di un database
relazionale. La sostituzione di una dipendenza deve essere esplicita e
non attivata per errore in produzione.

\section{5.7 Dati e integrazione con i sistemi
esistenti}\label{dati-e-integrazione-con-i-sistemi-esistenti}

Il database è spesso il vincolo più forte di una migrazione.
L'applicazione nuova può essere distribuita in poche ore; i dati possono
essere condivisi con altri servizi, avere anni di storia e contenere
anomalie che nessun modello documenta. Pianificare insieme schema,
codice, ownership, sincronizzazione e rollback.

\subsection{5.7.1 EF Core e scelta dell'accesso ai
dati}\label{ef-core-e-scelta-dellaccesso-ai-dati}

Entity Framework Core è un ORM che traduce query LINQ in operazioni sul
database e tiene traccia delle entità caricate. È adatto a molti casi
applicativi, ma non è obbligatorio per ogni query. SQL esplicito o un
micro-ORM possono essere più chiari per report, query complesse o
requisiti di prestazione. La scelta va fatta per workload e misurazioni,
non per fedeltà all'ORM sorgente.

{\def\LTcaptype{none} % do not increment counter
\begin{longtable}[]{@{}
  >{\raggedright\arraybackslash}p{(\linewidth - 4\tabcolsep) * \real{0.3333}}
  >{\raggedright\arraybackslash}p{(\linewidth - 4\tabcolsep) * \real{0.3333}}
  >{\raggedright\arraybackslash}p{(\linewidth - 4\tabcolsep) * \real{0.3333}}@{}}
\toprule\noalign{}
\begin{minipage}[b]{\linewidth}\raggedright
Concetto di altri ORM
\end{minipage} & \begin{minipage}[b]{\linewidth}\raggedright
EF Core
\end{minipage} & \begin{minipage}[b]{\linewidth}\raggedright
Verifica necessaria
\end{minipage} \\
\midrule\noalign{}
\endhead
\bottomrule\noalign{}
\endlastfoot
Session/entity manager & \texttt{DbContext} & Scope, tracking e
concorrenza \\
Entity/model & Entità CLR & Identità, navigazioni e mapping \\
Schema migration & EF migrations & SQL generato e deploy sicuro \\
Lazy loading & Opzionale, non predefinito in ogni scenario & Query
implicite e N+1 \\
Unit of work & \texttt{SaveChanges}/transazione & Confini e atomicità \\
\end{longtable}
}

\texttt{DbContext} è un'unità di lavoro breve, normalmente scoped per
richiesta. Non è thread-safe: non lanciare query concorrenti sullo
stesso contesto con \texttt{Task.WhenAll}. Usare query proiettate per
leggere solo i campi necessari e \texttt{AsNoTracking} quando il
risultato non verrà aggiornato tramite il context.

\begin{Shaded}
\begin{Highlighting}[]
\DataTypeTok{var}\NormalTok{ products }\OperatorTok{=}\NormalTok{ await db}\OperatorTok{.}\FunctionTok{Products}
    \OperatorTok{.}\FunctionTok{AsNoTracking}\OperatorTok{()}
    \OperatorTok{.}\FunctionTok{Where}\OperatorTok{(}\NormalTok{p }\OperatorTok{=\textgreater{}}\NormalTok{ p}\OperatorTok{.}\FunctionTok{IsActive}\OperatorTok{)}
    \OperatorTok{.}\FunctionTok{OrderBy}\OperatorTok{(}\NormalTok{p }\OperatorTok{=\textgreater{}}\NormalTok{ p}\OperatorTok{.}\FunctionTok{Name}\OperatorTok{)}
    \OperatorTok{.}\FunctionTok{Select}\OperatorTok{(}\NormalTok{p }\OperatorTok{=\textgreater{}} \KeywordTok{new} \FunctionTok{ProductListItem}\OperatorTok{(}
\NormalTok{        p}\OperatorTok{.}\FunctionTok{Id}\OperatorTok{,}\NormalTok{ p}\OperatorTok{.}\FunctionTok{Name}\OperatorTok{,}\NormalTok{ p}\OperatorTok{.}\FunctionTok{Price}\OperatorTok{,}\NormalTok{ p}\OperatorTok{.}\FunctionTok{Currency}\OperatorTok{))}
    \OperatorTok{.}\FunctionTok{ToListAsync}\OperatorTok{(}\NormalTok{cancellationToken}\OperatorTok{);}
\end{Highlighting}
\end{Shaded}

La proiezione esplicita rende visibile il contratto della query e riduce
trasferimento di dati. Non enumerare query infinite o materializzare
grandi tabelle senza paginazione. \texttt{IQueryable} esegue lavoro
quando enumerato; la query va mantenuta nel livello che conosce il
database per evitare sorprese.

\subsection{5.7.2 Modello, schema e
confini}\label{modello-schema-e-confini}

Non assumere che la struttura del modello sorgente si trasferisca
direttamente. Identificare chiavi primarie e naturali, vincoli,
nullability, precisione dei numeri, timezone, collations, indici,
relazioni e politiche di cancellazione. Confrontare il tipo logico e il
tipo fisico: \texttt{decimal} può avere scale diverse, un timestamp può
essere locale oppure UTC, un \texttt{string} può avere limiti e
collation.

Separare modello persistito e DTO quando la forma pubblica o le regole
del dominio differiscono dalla tabella. Non usare il database come unico
documento del dominio, ma non duplicare ogni campo in modelli senza un
confine utile. Le regole di integrità devono essere applicate anche nel
database per proteggere da race condition e percorsi non HTTP.

Prima di introdurre un nuovo database, capire quali altri componenti
leggono o scrivono lo schema attuale. Uno schema condiviso è un
contratto implicito. Migrare un servizio lasciando integrazioni non
documentate può interrompere sistemi che non fanno parte del repository.

\subsection{5.7.3 Migrazioni dello
schema}\label{migrazioni-dello-schema}

Le migrazioni EF Core descrivono evoluzioni dello schema; non
sostituiscono revisione e piano di deploy. Ispezionare lo SQL generato,
capire lock e durata, e distinguere modifiche additive da modifiche
distruttive. Un \texttt{ALTER\ TABLE} può bloccare una tabella grande;
il comportamento dipende dal provider e dalla versione del database.

Un pattern compatibile è \textbf{expand--migrate--contract}:

\begin{enumerate}
\def\labelenumi{\arabic{enumi}.}
\tightlist
\item
  aggiungere colonna o tabella compatibile con il codice precedente;
\item
  distribuire codice che legge/scrive nel formato compatibile;
\item
  copiare o trasformare i dati a batch e verificare conteggi;
\item
  passare la lettura al nuovo campo e osservare;
\item
  rimuovere il vecchio campo solo quando nessun consumer o rollback lo
  richiede.
\end{enumerate}

Non eseguire automaticamente migrazioni potenzialmente lunghe all'avvio
di ogni replica. Le istanze possono partire insieme e competere per
lock; il processo applicativo potrebbe non avere permessi DDL in
produzione. Preferire un passaggio controllato in pipeline o un job di
deployment con un solo proprietario e log chiari.

Per un rollback, la migrazione inversa deve essere realmente possibile.
Se il nuovo codice scrive valori non rappresentabili nel vecchio schema,
il downgrade non è sicuro. Conservare backup testati e definire una
strategia di ripristino coerente con RPO e RTO; un backup mai
ripristinato non è evidenza sufficiente.

\subsection{5.7.4 Migrare i dati}\label{migrare-i-dati}

Il piano dipende da volume, downtime consentito e possibilità di
mantenere sincronizzati i sistemi. Opzioni comuni:

\begin{itemize}
\tightlist
\item
  \textbf{Finestra di fermo:} bloccare scritture, esportare,
  trasformare, importare, verificare e riaprire. Semplice, ma richiede
  una finestra accettabile.
\item
  \textbf{Snapshot e delta:} copiare una fotografia iniziale, poi
  applicare le modifiche intervenute fino al cutover. Richiede un
  meccanismo affidabile per identificare il delta.
\item
  \textbf{Replica/CDC:} acquisire cambiamenti dal log del database o da
  un processo dedicato. Riduce il fermo, ma aggiunge dipendenze e
  necessità di controllo della consistenza.
\item
  \textbf{Migrazione per partizioni/tenant:} trasferire porzioni
  indipendenti e cambiare ownership progressivamente. Funziona solo se
  il confine è stabile.
\end{itemize}

Per ciascuna opzione descrivere il punto di consistenza, il
comportamento delle scritture durante la copia, l'ordine delle
trasformazioni e come riprendere dopo un'interruzione. Gli script devono
essere idempotenti o avere checkpoint persistenti. Registrare errori per
record e consentire la ripresa senza duplicare dati.

\subsection{5.7.5 Qualità e semantica dei
dati}\label{qualituxe0-e-semantica-dei-dati}

I dati storici possono violare le regole attuali: stringhe vuote al
posto di null, date impossibili, duplicati, codici fuori dominio, valute
miste o relazioni orfane. Non correggere silenziosamente senza una
regola approvata. Classificare gli elementi in validi, trasformabili, da
mettere in quarantena e bloccanti.

Per date e orari distinguere istante UTC, orario locale e data civile.
\texttt{DateTime} senza \texttt{Kind} esplicito può essere ambiguo;
\texttt{DateTimeOffset} preserva offset ma non identifica da solo il
fuso orario nominato. Convertire secondo il significato di business, non
aggiungendo semplicemente \texttt{Z} alle stringhe.

Per denaro definire valuta, scala e arrotondamento; non usare floating
point binario per valori monetari senza una motivazione. Per testo
considerare Unicode, normalizzazione e collation: due stringhe
visivamente identiche possono avere rappresentazioni differenti. Per
identificativi preservare zeri iniziali se fanno parte del valore.

\subsection{5.7.6 Riconciliazione
verificabile}\label{riconciliazione-verificabile}

Il confronto tra origine e target deve andare oltre il conteggio righe.
Verificare:

\begin{itemize}
\tightlist
\item
  conteggi per tabella, tenant e intervallo temporale;
\item
  checksum o hash di colonne canoniche;
\item
  riferimenti e vincoli;
\item
  aggregati di business, per esempio totali per stato;
\item
  campioni mirati di record complessi;
\item
  dati modificati durante il delta e record scartati.
\end{itemize}

Definire tolleranze solo dove la trasformazione le giustifica, come
precisione numerica o normalizzazione. Conservare report firmati o
immutabili del processo, con timestamp, versione dello script e
conteggi. Un esito ``completato'' senza numeri è difficile da auditare.

\subsection{5.7.7 Transazioni, concorrenza e
query}\label{transazioni-concorrenza-e-query}

Una transazione database protegge un insieme di modifiche all'interno
del provider e del suo contesto; non rende atomica una chiamata HTTP più
una scrittura. Non tenere una transazione aperta mentre si attende un
servizio remoto. Per coordinare DB ed eventi usare pattern come outbox
con consumer idempotente.

Gestire concorrenza ottimistica con token o timestamp e restituire un
conflitto comprensibile quando una versione è cambiata. Testare query
reali con il provider di produzione: SQLite in memoria può differire per
tipi, collations, transazioni e funzioni. I test rapidi con provider
sostitutivo sono utili ma non sostituiscono test d'integrazione sul
motore effettivo.

Misurare query e piani d'esecuzione prima di aggiungere repository o
cache. Attenzione a N+1, caricamento eccessivo, assenza di indici e
paginazione con offset su dataset enormi. Non interpolare SQL con input;
usare parametri e permessi minimi.

\subsection{5.7.8 Connessioni e
continuità}\label{connessioni-e-continuituxe0}

Le stringhe di connessione provengono da configurazione protetta.
Impostare pool, timeout, TLS e retry secondo provider; una connessione
esaurita può apparire come lentezza dell'API. Il pool riduce l'overhead
ma non aumenta la capacità del database: moltiplicare repliche
moltiplica anche le connessioni potenziali.

Durante la convivenza, definire quale sistema può scrivere ciascuna
entità e come viene mantenuta la consistenza. Evitare di far dipendere
la nuova API direttamente da uno schema che il vecchio sistema può
cambiare senza coordinamento. Un contratto dati versionato o un adapter
temporaneo può isolare la migrazione e rendere esplicita la dipendenza
da rimuovere.

\section{5.8 Sicurezza e identità}\label{sicurezza-e-identituxe0}

La sicurezza non è un passaggio finale di hardening. Cambiare framework
e identità modifica cookie, token, chiavi, policy e trust boundary. Una
migrazione può mantenere gli endpoint e allo stesso tempo disconnettere
gli utenti, accettare token non validi o rivelare dati. Definire i
requisiti e verificare il flusso end-to-end prima del cutover.

\subsection{5.8.1 Iniziare da threat model e
inventario}\label{iniziare-da-threat-model-e-inventario}

Identificare asset (dati personali, credenziali, pagamenti, sessioni),
attori, confini di fiducia e minacce rilevanti. Per ogni flusso
chiedersi chi può invocarlo, come viene autenticato, quali permessi
servono, quali input controlla e cosa accade se la dipendenza fallisce.

Mappare:

\begin{itemize}
\tightlist
\item
  provider di identità e modalità di login;
\item
  cookie, sessioni, token access e refresh, durata e revoca;
\item
  ruoli, claims, tenant e regole d'autorizzazione;
\item
  segreti applicativi, chiavi di firma e certificati;
\item
  dati sensibili nei log, backup, cache e telemetria;
\item
  protezioni esistenti, dipendenze vulnerabili e procedure d'incidente.
\end{itemize}

Le policy scritte nel codice possono essere incomplete rispetto alla
configurazione del gateway o del provider. Verificare la catena
completa, inclusi reverse proxy e client. Fare approvare i cambiamenti
da chi è responsabile di sicurezza e identità.

\subsection{5.8.2 Autenticazione e autorizzazione non sono la stessa
cosa}\label{autenticazione-e-autorizzazione-non-sono-la-stessa-cosa}

L'autenticazione stabilisce l'identità; l'autorizzazione decide cosa
quell'identità può fare. In ASP.NET Core la pipeline usa handler di
autenticazione e policy:

\begin{Shaded}
\begin{Highlighting}[]
\NormalTok{builder}\OperatorTok{.}\FunctionTok{Services}\OperatorTok{.}\FunctionTok{AddAuthentication}\OperatorTok{()}
    \OperatorTok{.}\FunctionTok{AddJwtBearer}\OperatorTok{();}

\NormalTok{builder}\OperatorTok{.}\FunctionTok{Services}\OperatorTok{.}\FunctionTok{AddAuthorization}\OperatorTok{(}\NormalTok{options }\OperatorTok{=\textgreater{}}
\OperatorTok{\{}
\NormalTok{    options}\OperatorTok{.}\FunctionTok{AddPolicy}\OperatorTok{(}\StringTok{"CatalogEditors"}\OperatorTok{,}\NormalTok{ policy }\OperatorTok{=\textgreater{}}
\NormalTok{        policy}\OperatorTok{.}\FunctionTok{RequireClaim}\OperatorTok{(}\StringTok{"permission"}\OperatorTok{,} \StringTok{"catalog.write"}\OperatorTok{));}
\OperatorTok{\});}

\NormalTok{app}\OperatorTok{.}\FunctionTok{UseAuthentication}\OperatorTok{();}
\NormalTok{app}\OperatorTok{.}\FunctionTok{UseAuthorization}\OperatorTok{();}

\NormalTok{app}\OperatorTok{.}\FunctionTok{MapPost}\OperatorTok{(}\StringTok{"/api/products"}\OperatorTok{,}\NormalTok{ CreateProduct}\OperatorTok{)}
   \OperatorTok{.}\FunctionTok{RequireAuthorization}\OperatorTok{(}\StringTok{"CatalogEditors"}\OperatorTok{);}
\end{Highlighting}
\end{Shaded}

Il frammento mostra la forma generale, non una configurazione JWT
completa: emittente (\texttt{iss}), destinatario (\texttt{aud}), chiavi,
HTTPS, scarto orario e validazione dei claim vanno configurati secondo
l'identity provider. Non disabilitare la validazione della firma o del
destinatario per far passare un test.

Ruoli e claims hanno semantiche diverse. Un ruolo descrive appartenenza
a una categoria; un claim può rappresentare un attributo o permesso. Non
affidarsi a un claim ricevuto da un client non autenticato. Per
multi-tenant, verificare che l'identificativo tenant sia vincolato
all'identità e che ogni query applichi il perimetro corretto.

\subsection{5.8.3 Migrare password e
account}\label{migrare-password-e-account}

Gli hash delle password non sono password reversibili. Se il sistema
nuovo usa un algoritmo o formato diverso, le opzioni includono:

\begin{itemize}
\tightlist
\item
  mantenere temporaneamente la verifica tramite il sistema identità
  esistente;
\item
  validare al primo login con il vecchio verificatore e aggiornare
  l'hash in modo sicuro;
\item
  richiedere reset password se la verifica compatibile non è
  disponibile;
\item
  migrare verso un identity provider condiviso con processo supportato.
\end{itemize}

Non abbassare il costo dell'hash o introdurre una password universale
per semplificare il passaggio. Non esportare hash in file non protetti.
Proteggere il flusso di reset, limitare tentativi e non rivelare se un
account esiste. Pianificare account bloccati, MFA, utenti disabilitati e
audit.

\subsection{5.8.4 Cookie, token e continuità di
sessione}\label{cookie-token-e-continuituxe0-di-sessione}

I cookie devono avere \texttt{Secure}, \texttt{HttpOnly} e una policy
\texttt{SameSite} coerente con i flussi di login. Se contengono
credenziali, gestire CSRF per richieste che modificano stato. La
protezione CSRF non è sostituita da CORS.

JWT è un formato di token, non una strategia di autorizzazione completa.
Validare firma, emittente (\texttt{iss}), destinatario (\texttt{aud}),
scadenza e algoritmo consentito; limitare durata e proteggere refresh
token. Non usare token come sessione revocabile senza un meccanismo di
revoca o rotazione. Ridurre i dati sensibili contenuti nel payload: il
payload firmato non è necessariamente cifrato.

I cookie di ASP.NET Core Identity o la Data Protection API possono
essere specifici dell'applicazione e delle chiavi usate. Se durante il
rollout vecchio e nuovo sistema devono accettare lo stesso cookie,
verificare esplicitamente formato, scopo, chiavi e compatibilità
supportata. Non improvvisare la condivisione delle chiavi:
comprometterle può consentire la forgiatura di cookie o la decifratura
di dati protetti.

\subsection{5.8.5 Data Protection API e
chiavi}\label{data-protection-api-e-chiavi}

ASP.NET Core Data Protection protegge payload applicativi, inclusi
cookie e token antiforgery. In sviluppo le chiavi possono essere locali;
in produzione le istanze devono avere una configurazione durevole e
coerente quando devono decifrare gli stessi payload. Conservare le
chiavi in storage protetto, limitarne l'accesso, pianificare rotazione e
backup e impostare un application name appropriato per isolare
applicazioni.

Un container effimero non deve essere l'unico luogo dove risiedono
chiavi che servono dopo il riavvio. D'altra parte, condividere
indiscriminatamente lo stesso key ring fra applicazioni aumenta
l'impatto di una compromissione. Testare restart, scale-out e rotazione
in un ambiente simile alla produzione.

\subsection{5.8.6 Segreti e configurazione
sicura}\label{segreti-e-configurazione-sicura}

Usare identità gestite o secret store quando disponibili, permessi
minimi, rotazione e audit. Evitare chiavi in
\texttt{appsettings.Production.json}, Dockerfile, immagini, repository o
log CI. Un secret presente nella history Git va considerato compromesso:
rimuoverlo dal file non revoca la credenziale.

Separare segreti e configurazione ordinaria, ma non trattare un URL
interno o un nome utente come sicuro per definizione. Tutto ciò che
compare nelle eccezioni o negli header può finire in un log. Definire
redazione e classificazione dei dati per log, tracing e metriche.

\subsection{5.8.7 Difese per API web}\label{difese-per-api-web}

\begin{itemize}
\tightlist
\item
  Validare input sul server, limitare dimensione e complessità delle
  richieste.
\item
  Usare query parametrizzate; non costruire SQL concatenando input.
\item
  Codificare output secondo il contesto e non fidarsi di dati già
  salvati.
\item
  Configurare CORS con origini esplicite; non è un controllo di accesso
  per client non-browser.
\item
  Impostare TLS e forwarded headers solo per proxy fidati.
\item
  Applicare rate limiting, timeout e limiti di concorrenza appropriati.
\item
  Ridurre permessi di database, file system e identità cloud.
\item
  Mantenere pacchetti e immagini aggiornati e scansionare le dipendenze.
\end{itemize}

I middleware di sicurezza devono essere ordinati e testati in hosting
reale. Una policy CORS corretta in locale può essere diversa dietro
gateway. Non esporre stack trace in produzione; associare l'errore a un
correlation id che il supporto può cercare.

\subsection{5.8.8 Test di sicurezza nella
transizione}\label{test-di-sicurezza-nella-transizione}

Provare almeno: anonimo su route protetta; utente autenticato senza
permesso; permesso valido; token scaduto, malfirmato o con destinatario
(\texttt{aud}) errato; tenant diverso; cookie senza protezione CSRF;
credenziale revocata; limite di rate; input e payload eccessivi.
Verificare che gli errori non rivelino dati e che i log non contengano
token o password.

Mantenere una matrice delle policy per endpoint. Ogni route deve avere
una decisione esplicita: anonima, autenticata o protetta da policy
nominata. Le route nuove non devono diventare anonime per omissione.
Integrare test automatici e una verifica di configurazione in CI.

\subsection{5.8.9 Cutover dell'identità}\label{cutover-dellidentituxe0}

Durante una migrazione incrementale è possibile che alcuni endpoint
usino il vecchio provider e altri il nuovo. Definire chi emette il
token, chi lo valida e quali claims sono autorevoli. Se si usa un
gateway per tradurre identità, documentare i trust boundary e impedire
che il client possa falsificare gli header interni.

Preparare una strategia per sessioni non trasferibili: riautenticazione
comunicata agli utenti, durata della transizione, revoca delle sessioni
e supporto. Misurare i login falliti dopo il cutover e predisporre un
rollback che non riattivi credenziali compromesse. La continuità non
giustifica mantenere per sempre un algoritmo o un cookie legacy
insicuro.

\section{5.9 Testing e qualità}\label{testing-e-qualituxe0}

La migrazione è l'occasione per rendere esplicito il comportamento del
sistema, non per sostituire i test del vecchio stack con test che
verificano solo dettagli interni del nuovo. I test più preziosi
descrivono ciò che client e utenti possono osservare, e permettono di
confrontare le due implementazioni sugli stessi casi.

\subsection{5.9.1 Test di
caratterizzazione}\label{test-di-caratterizzazione}

Un test di caratterizzazione registra il comportamento attuale, anche
quando la documentazione è incompleta. Si preparano input controllati,
si invocano route o funzioni del sistema sorgente e si salvano status,
header e payload normalizzati. Questi test diventano una specifica di
compatibilità per il target.

La normalizzazione va limitata a campi realmente variabili, come
timestamp generati o identificativi casuali. Normalizzare interi payload
può nascondere regressioni. Per dati sensibili usare fixture sintetiche
o anonimizzate; non importare dump di produzione nei test senza
autorizzazione e protezioni.

Per endpoint rappresentativi raccogliere:

\begin{itemize}
\tightlist
\item
  risposta valida e codici di successo;
\item
  input mancante, malformato e ai limiti;
\item
  risorsa assente e conflitto;
\item
  chiamante anonimo e permessi insufficienti;
\item
  timeout o indisponibilità di una dipendenza;
\item
  richieste duplicate e concorrenza.
\end{itemize}

Il test di caratterizzazione non dichiara che il comportamento corrente
sia ideale. Un bug noto può essere registrato come comportamento da
cambiare, con un test che viene aggiornato nella stessa modifica della
decisione e con l'approvazione del proprietario del contratto.

\subsection{5.9.2 Piramide e livelli di
test}\label{piramide-e-livelli-di-test}

\textbf{Unit test} verificano funzioni e regole isolate, rapidamente e
senza rete. Sono adatti a validazioni, calcoli, policy e transizioni di
stato. Non devono ricreare l'intero framework con mock che duplicano
l'implementazione.

\textbf{Test d'integrazione} esercitano più componenti: hosting, route,
serializzazione, DI, database e adapter. In ASP.NET Core
\texttt{WebApplicationFactory\textless{}TEntryPoint\textgreater{}} avvia
l'applicazione nel processo di test e permette richieste HTTP in
memoria. Sostituire configurazione o dipendenze in modo controllato;
verificare anche la pipeline vera.

\textbf{Test end-to-end} attraversano deployment, gateway e servizi
integrati. Sono più lenti e sensibili all'ambiente; riservarli ai flussi
business critici, mantenere dati isolati e avere pulizia affidabile.

La piramide è una guida alla velocità e al feedback, non una quota
rigida. Un test d'integrazione diretto sul database può essere più
attendibile di un grande insieme di mock che non simula transazioni e
SQL.

\subsection{5.9.3 Test HTTP con
WebApplicationFactory}\label{test-http-con-webapplicationfactory}

Un test concettuale può verificare un contratto HTTP:

\begin{Shaded}
\begin{Highlighting}[]
\KeywordTok{public} \KeywordTok{sealed} \KeywordTok{class} \FunctionTok{ProductApiTests}\OperatorTok{(}
\NormalTok{    WebApplicationFactory}\OperatorTok{\textless{}}\NormalTok{Program}\OperatorTok{\textgreater{}}\NormalTok{ factory}\OperatorTok{)}
    \OperatorTok{:}\NormalTok{ IClassFixture}\OperatorTok{\textless{}}\NormalTok{WebApplicationFactory}\OperatorTok{\textless{}}\NormalTok{Program}\OperatorTok{\textgreater{}\textgreater{}}
\OperatorTok{\{}
    \OperatorTok{[}\NormalTok{Fact}\OperatorTok{]}
    \KeywordTok{public}\NormalTok{ async Task }\FunctionTok{Missing\_product\_returns\_not\_found}\OperatorTok{()}
    \OperatorTok{\{}
        \KeywordTok{using} \DataTypeTok{var}\NormalTok{ client }\OperatorTok{=}\NormalTok{ factory}\OperatorTok{.}\FunctionTok{CreateClient}\OperatorTok{();}
        \KeywordTok{using} \DataTypeTok{var}\NormalTok{ response }\OperatorTok{=}
\NormalTok{            await client}\OperatorTok{.}\FunctionTok{GetAsync}\OperatorTok{(}\StringTok{"/api/products/987654"}\OperatorTok{);}

\NormalTok{        Assert}\OperatorTok{.}\FunctionTok{Equal}\OperatorTok{(}\NormalTok{HttpStatusCode}\OperatorTok{.}\FunctionTok{NotFound}\OperatorTok{,}\NormalTok{ response}\OperatorTok{.}\FunctionTok{StatusCode}\OperatorTok{);}
    \OperatorTok{\}}
\OperatorTok{\}}
\end{Highlighting}
\end{Shaded}

Il progetto deve esporre il punto d'ingresso al test se necessario e
configurare un database isolato. Non riutilizzare un database condiviso
tra esecuzioni concorrenti senza namespace o transazioni pulite. Per
test di persistenza, verificare vincoli, migrazioni e comportamento del
provider adottato in produzione.

Testare le risposte come farebbe un client: status, content type,
header, schema e valori. Non verificare il nome di una classe interna
quando l'obiettivo è la compatibilità API. Test specifici di unità
possono invece verificare il dominio senza attraversare HTTP.

\subsection{5.9.4 Parità differenziale}\label{parituxe0-differenziale}

Quando entrambi i sistemi sono avviabili, inviare lo stesso caso a
sorgente e target e confrontare gli output. Il test differenziale può
essere automatizzato su una collezione di casi deterministici. Escludere
o normalizzare solo campi non deterministici, mantenendo visibili
differenze nei codici, nei valori e negli errori.

Non inviare automaticamente richieste di scrittura reali a entrambe le
versioni: potrebbero creare effetti duplicati, inviare email o
addebitare un pagamento. Per i casi mutativi usare database isolati,
simulazioni senza effetti collaterali o una modalità di shadow che
calcola la risposta senza committare.

Il confronto deve considerare anche proprietà mancanti e ordine delle
liste quando l'ordine è contrattuale. Per importi, date e testo,
confrontare con regole esplicite; non applicare una tolleranza generica
che renda irrilevanti differenze di business.

\subsection{5.9.5 Dipendenze nei test}\label{dipendenze-nei-test}

Mock e fake hanno valore quando isolano un confine, per esempio un
servizio di pagamento. Non mockare il database con una collezione in
memoria e concludere che le transazioni reali funzionano. Usare
container effimeri o istanze dedicate per testare provider SQL, code o
identity server quando la compatibilità è importante.

Ogni fake deve avere semantica dichiarata. Un fake cache senza scadenza
non verifica TTL; un fake storage senza errori non prova retry; un fake
repository che restituisce sempre l'entità non esercita i vincoli.
Preferire test a contratto per verificare che adapter diversi rispettino
lo stesso protocollo.

\subsection{5.9.6 Flakiness, dati e
parallelismo}\label{flakiness-dati-e-parallelismo}

Un test che dipende dall'orologio reale, da un endpoint esterno o da
dati condivisi può essere intermittente. Iniettare un clock quando il
tempo è parte del comportamento, controllare il seed e isolare risorse.
Non disabilitare l'esecuzione parallela dell'intera suite per nascondere
una race: identificare lo stato condiviso e isolare il test interessato.

I test devono poter essere ripetuti. Ogni esecuzione dovrebbe creare e
rimuovere i propri dati, usare nomi univoci e produrre log utili in caso
di fallimento. Non riutilizzare credenziali di produzione. La cleanup
deve funzionare anche dopo un errore, senza cancellare database o
namespace condivisi.

\subsection{5.9.7 Qualità e pipeline CI}\label{qualituxe0-e-pipeline-ci}

Una pipeline minima esegue restore, build Release, test, analisi e
pubblicazione di artefatti. Aggiungere scansione dipendenze,
formattazione e analisi statiche secondo il rischio e gli standard del
team. Bloccare la release su fallimenti significativi e rendere i
risultati rintracciabili con commit, versione SDK e configurazione.

\begin{Shaded}
\begin{Highlighting}[]
\NormalTok{dotnet restore}
\NormalTok{dotnet format {-}{-}verify{-}no{-}changes}
\NormalTok{dotnet build {-}c Release {-}{-}no{-}restore}
\NormalTok{dotnet test {-}c Release {-}{-}no{-}build {-}{-}logger "trx"}
\NormalTok{dotnet publish src/Catalogo.Api {-}c Release {-}o artifacts/publish}
\end{Highlighting}
\end{Shaded}

La pipeline deve funzionare da checkout pulito e non dipendere da file
generati presenti sulla workstation. Separare test che richiedono
servizi esterni e rendere espliciti prerequisiti e cleanup. Verificare
che il pacchetto pubblicato sia lo stesso che verrà promosso, non
ricostruire da sorgenti diversi in ogni ambiente.

\subsection{5.9.8 Copertura orientata al
rischio}\label{copertura-orientata-al-rischio}

La percentuale di line coverage non misura quanto i test proteggono il
contratto. Dare priorità a flussi con impatto alto, transizioni di
stato, permessi, calcoli finanziari, migrazione dati e fallback. Un
singolo test di integrazione ben progettato può coprire un rischio non
espresso da decine di test di getter.

Per ogni capacità migrata, associare requisito, test, output e owner. I
criteri di uscita possono richiedere: test di parità superati;
regressioni spiegate; contratti aggiornati; test sul provider reale;
nessun difetto bloccante; report di qualità archiviato. Se una parte non
è testabile, registrare il rischio e una mitigazione alternativa.

\section{5.10 Prestazioni e comportamento del
runtime}\label{prestazioni-e-comportamento-del-runtime}

Una migrazione cambia runtime, librerie, scheduling, serializer e
modalità di accesso ai dati. Il nuovo servizio può risultare più
efficiente, equivalente o peggiore: la risposta si misura sul workload
reale. Ottimizzare sulla reputazione del linguaggio o su un benchmark
isolato porta a decisioni costose e poco attendibili.

\subsection{5.10.1 Baseline e carichi
confrontabili}\label{baseline-e-carichi-confrontabili}

Scegliere scenari che rappresentino uso, picchi e operazioni costose:
lettura cache hit, miss, query DB, payload grande, dipendenza lenta e
endpoint mutativo. Usare dataset e distribuzione dei dati confrontabili;
una tabella con mille righe non simula una con milioni.

Misurare latenza p50/p95/p99, throughput, errori, CPU, memoria,
allocazioni, connessioni e tempi di avvio. Includere warm-up e durata
sufficiente a osservare garbage collection e saturazione delle
dipendenze. Separare misure del client e del server: una rete più lenta
non è una regressione del codice applicativo.

Per ogni test annotare versione runtime, build Release, hardware,
container limits, numero istanze, configurazione, seed e commit. I
confronti devono essere ripetibili; se il rumore è alto, raccogliere più
campioni e indicare intervalli, non il singolo valore migliore.

\subsection{5.10.2 Asincronia e
concorrenza}\label{asincronia-e-concorrenza}

ASP.NET Core gestisce molte richieste concorrenti senza assegnare
necessariamente un thread dedicato a ogni attesa I/O.
\texttt{async}/\texttt{await} consente di restituire il thread mentre si
aspetta rete o disco. Non rende un calcolo CPU-bound più rapido, né
elimina i limiti di database e upstream.

Evitare \texttt{.Result}, \texttt{.Wait()} e
\texttt{GetAwaiter().GetResult()} sul percorso HTTP: il blocco occupa
thread e sotto carico può causare starvation. Propagare
\texttt{CancellationToken} fino ai client e alla persistenza. Se il
client annulla la richiesta, il server dovrebbe interrompere lavoro
inutile quando è sicuro farlo.

Non avviare operazioni parallele sullo stesso \texttt{DbContext}. Se più
query indipendenti sono davvero necessarie, usare scope/contesti
distinti con limite di concorrenza e verificare che il database possa
sostenere il carico. Il parallelismo indiscriminato peggiora la latenza
quando satura pool e connessioni.

\subsection{5.10.3 Memoria e garbage
collection}\label{memoria-e-garbage-collection}

Il runtime .NET usa garbage collector generazionale; allocazioni
temporanee sono spesso economiche, ma allocazioni eccessive aumentano
pressione GC e latenza. Profilare prima di introdurre pooling o
strutture più complesse. Riutilizzare buffer in modo scorretto può
causare corruzione o esposizione di dati tra richieste.

Evitare di leggere file o payload illimitati in memoria. Applicare
limiti, streaming e backpressure dove appropriato. Cache, singleton e
collezioni statiche possono trattenere oggetti per tutta la vita del
processo; monitorare dimensioni e politiche di eviction. Un processo
containerizzato deve rispettare limiti memoria: essere stabile sulla
workstation non prova che lo sia con un memory limit ridotto.

\subsection{5.10.4 Query, serializzazione e
caching}\label{query-serializzazione-e-caching}

La causa dominante di lentezza è spesso I/O o query non efficienti.
Usare log delle query e strumenti del database per diagnosticare N+1,
full scan, lock e attese. Proiettare i campi necessari, paginare,
aggiungere indici motivati dal piano e misurare dopo la modifica.

La serializzazione JSON può diventare significativa per payload grandi o
grafi complessi. Definire DTO, evitare cicli nelle navigazioni e
misurare allocazioni e dimensione della risposta. Cambiare serializer o
opzioni di naming può rompere il contratto, anche se migliora le
prestazioni.

La cache è utile quando riduce lavoro ripetuto e la staleness è
accettabile. Definire TTL, invalidazione, chiave, isolamento tenant,
limite memoria e comportamento in caso di errore. Misurare hit ratio e
costo di riempimento. Non aggiungere cache per compensare query senza
aver compreso la causa.

\subsection{5.10.5 Profiling e Native AOT}\label{profiling-e-native-aot}

Usare profiler e tracing per individuare il tratto costoso, confrontando
CPU sample, allocation, thread pool e dipendenze. Un flame graph aiuta a
trovare hot path, ma va letto insieme a metriche applicative. Gli
strumenti di tracing devono essere configurati per minimizzare overhead
e non registrare payload sensibili.

Trimming e Native AOT possono ridurre footprint o tempo di avvio per
scenari specifici, ma imposero vincoli a reflection, serializzazione,
plugin e librerie. Abilitare l'analisi e correggere warning, non
silenziarli indiscriminatamente. Misurare cold start, dimensione e
throughput nel target effettivo. Non assumere che AOT migliori un
servizio dominato da I/O o database.

\subsection{5.10.6 Prestazioni come criterio di
cutover}\label{prestazioni-come-criterio-di-cutover}

Definire budget per operazione e soglie di rollback prima del rollout.
Durante canary confrontare gruppi equivalenti per traffico, carico e
dipendenze. Osservare code, errori e saturazione, non soltanto latenza
media. Se il target peggiora, poterlo disattivare deve essere più rapido
che applicare una patch d'emergenza.

La decisione finale non è ``più veloce a ogni costo''. Un sistema con
latenza entro budget, meno incidenti e costo prevedibile può essere
preferibile a un microbenchmark migliore ma a scapito di complessità e
manutenibilità.

\section{5.11 Container, CI/CD e
produzione}\label{container-cicd-e-produzione}

La migrazione termina quando il sistema è esercibile dal team che lo
mantiene: build riproducibile, artifact identificabile, configurazione
esterna, deploy osservabile, rollback eseguibile e responsabilità di
supporto chiare. Un endpoint che funziona sul computer dello
sviluppatore non è ancora un servizio migrato.

\subsection{5.11.1 Pubblicazione e
artefatti}\label{pubblicazione-e-artefatti}

\texttt{dotnet\ publish\ -c\ Release} prepara output per un runtime e un
ambiente di distribuzione. La pipeline deve compilare una volta e
promuovere lo stesso artefatto tra gli ambienti, aggiungendo
configurazione esterna. Ricostruire a ogni stage può produrre binari
differenti senza che il team se ne accorga.

Registrare commit, versione SDK, dipendenze e provenance dell'immagine.
Per applicazioni framework-dependent, il server deve avere un runtime
compatibile; per self-contained, il pacchetto include runtime e richiede
patching e dimensioni gestite diversamente. La scelta va documentata in
relazione a hosting e policy di aggiornamento.

\subsection{5.11.2 Dockerfile multi-stage}\label{dockerfile-multi-stage}

Un Dockerfile multi-stage separa compilazione e runtime:

\begin{Shaded}
\begin{Highlighting}[]
\KeywordTok{FROM}\NormalTok{ mcr.microsoft.com/dotnet/sdk:10.0 }\KeywordTok{AS}\NormalTok{ build}
\KeywordTok{WORKDIR}\NormalTok{ /src}
\KeywordTok{COPY}\NormalTok{ [}\StringTok{"src/Catalogo.Api/Catalogo.Api.csproj"}\NormalTok{, }\StringTok{"src/Catalogo.Api/"}\NormalTok{]}
\KeywordTok{COPY}\NormalTok{ [}\StringTok{"src/Catalogo.Infrastructure/Catalogo.Infrastructure.csproj"}\NormalTok{, }\StringTok{"src/Catalogo.Infrastructure/"}\NormalTok{]}
\KeywordTok{RUN} \ExtensionTok{dotnet}\NormalTok{ restore }\StringTok{"src/Catalogo.Api/Catalogo.Api.csproj"}
\KeywordTok{COPY}\NormalTok{ . .}
\KeywordTok{RUN} \ExtensionTok{dotnet}\NormalTok{ publish }\StringTok{"src/Catalogo.Api/Catalogo.Api.csproj"} \DataTypeTok{\textbackslash{}}
    \AttributeTok{{-}c}\NormalTok{ Release }\AttributeTok{{-}o}\NormalTok{ /app/publish }\AttributeTok{{-}{-}no{-}restore}

\KeywordTok{FROM}\NormalTok{ mcr.microsoft.com/dotnet/aspnet:10.0 }\KeywordTok{AS}\NormalTok{ final}
\KeywordTok{WORKDIR}\NormalTok{ /app}
\KeywordTok{COPY} \OperatorTok{{-}{-}from=build}\NormalTok{ /app/publish .}
\KeywordTok{USER}\NormalTok{ $APP\_UID}
\KeywordTok{ENTRYPOINT}\NormalTok{ [}\StringTok{"dotnet"}\NormalTok{, }\StringTok{"Catalogo.Api.dll"}\NormalTok{]}
\end{Highlighting}
\end{Shaded}

Adattare path, identità e immagini alla versione supportata e alla
policy organizzativa. Non copiare segreti nel build context; usare
\texttt{.dockerignore} per escludere \texttt{.git}, artefatti e file
locali. Pinning e aggiornamento delle immagini devono essere governati:
fissare un digest aumenta riproducibilità, ma richiede un processo per
ricevere patch.

Eseguire con utente non root quando supportato; configurare porte e
variabili dal runtime, non codificarle nel Dockerfile. Analizzare
l'immagine, limitare capability e filesystem scrivibile, e verificare i
certificati CA necessari alle chiamate HTTPS. Un container piccolo non è
un container sicuro se contiene vulnerabilità non aggiornate.

\subsection{5.11.3 CI/CD}\label{cicd}

La pipeline dovrebbe effettuare restore, build, test, analisi,
pubblicazione, scansione, firma/provenance se prevista e deploy. I
passaggi devono fallire in modo esplicito. Non stampare token in output
e limitare i permessi del runner. Usare credenziali effimere o federate
quando disponibili.

Separare build da deploy: il job di build produce artefatto immutabile;
gli stage successivi applicano configurazione e promuovono lo stesso
digest. Le approvazioni manuali possono essere necessarie per ambienti
sensibili, ma non devono sostituire test automatizzati. Conservare log e
report in un sistema accessibile al team durante incidenti.

La strategia di database va coordinata con la pipeline: le migrazioni
non devono partire in parallelo da ogni replica. Le modifiche
incompatibili richiedono deployment in più fasi. La pipeline dovrebbe
rendere visibili script SQL generati e l'approvazione necessaria.

\subsection{5.11.4 Reverse proxy e
hosting}\label{reverse-proxy-e-hosting}

In produzione Kestrel può essere dietro IIS, Nginx, un ingress
Kubernetes o un gateway cloud. Configurare forwarded headers da proxy
fidati per ricostruire schema e IP originali. Se si accettano header
forwarded da qualunque client, un chiamante può falsificare indirizzi o
aggirare regole basate sull'IP.

Terminare TLS nel punto previsto, proteggere il collegamento tra proxy e
applicazione quando richiesto e configurare redirect HTTPS coerenti.
Verificare \texttt{PathBase} se il gateway pubblica l'app sotto un
prefisso. Controllare upload size, timeout, buffering e idle timeout a
ogni livello: il limite effettivo è spesso quello più basso della
catena.

\subsection{5.11.5 Health, osservabilità e
shutdown}\label{health-osservabilituxe0-e-shutdown}

Distinguere:

\begin{itemize}
\tightlist
\item
  \textbf{liveness:} il processo è in grado di continuare a funzionare?
\item
  \textbf{readiness:} l'istanza può ricevere traffico?
\item
  \textbf{startup:} l'inizializzazione è completata?
\end{itemize}

Un health check che prova ogni dipendenza a ogni intervallo può
amplificare un guasto. Separare controlli superficiali da diagnostica
profonda e usare timeout. Non esporre endpoint di diagnostica
dettagliati al pubblico.

Usare log strutturati con trace/correlation id, metriche aggregate e
distributed tracing per seguire chiamate tra servizi. Evitare
identificativi personali o token nei label delle metriche: cardinalità e
privacy. Definire dashboard e alert collegati a obiettivi di servizio,
con runbook e owner.

Lo shutdown deve cessare l'accettazione di nuovo traffico e consentire
il completamento del lavoro entro il grace period dell'orchestratore. I
worker devono ricevere cancellation token e fermarsi in modo
controllato. Testare terminazione con richieste in corso e drain del
load balancer.

\subsection{5.11.6 Scalabilità, storage e
resilienza}\label{scalabilituxe0-storage-e-resilienza}

Un'applicazione web stateless è più facile da scalare, ma cache locale,
sessione, chiavi Data Protection, file caricati e job in memoria
introducono stato di istanza. Documentare dove vive ciascun dato e cosa
succede se un nodo viene rimosso.

Usare storage condiviso per file che devono sopravvivere al container;
non considerare il filesystem locale una fonte durevole. Se sono
necessari sticky session, documentare i limiti e pianificare una
migrazione verso stato distribuito quando opportuno. Limitare pool di
connessioni in relazione al numero massimo di repliche e alla capacità
database.

Resilienza non significa retry illimitato. Configurare timeout, retry
limitati con backoff, circuit breaker o bulkhead secondo il failure
mode. L'operazione deve avere semantica idempotente o un identificatore
di deduplicazione. Un upstream degradato non deve esaurire thread,
connessioni e code del servizio migrato.

\subsection{5.11.7 Rollout progressivo e rollback
operativo}\label{rollout-progressivo-e-rollback-operativo}

Con canary o blue/green, spostare percentuali di traffico e confrontare
errori e latenza su gruppi simili. Impostare soglie automatiche e un
responsabile umano per l'arresto. Un rollout graduale senza metriche
segmentate è solo un deploy lento.

Il rollback del codice deve essere compatibile con lo schema dati e con
gli eventi già prodotti. Se i messaggi sono stati consumati o i dati
trasformati, tornare al binario precedente potrebbe non bastare.
Preparare procedure per fermare consumer, drenare code, ripristinare
dati o completare una compensazione.

Provare l'intero runbook in un ambiente rappresentativo: identificare la
versione attiva, cambiare routing, confermare i segnali, annunciare lo
stato e raccogliere evidenze. Registrare durata e impedimenti; assegnare
le correzioni prima del cutover reale.

\subsection{5.11.8 Costi, supporto e
dismissione}\label{costi-supporto-e-dismissione}

Calcolare costo complessivo includendo doppio runtime durante la
transizione, database, log, traffico, build runner, licenze e tempo
operativo. La coesistenza dovrebbe avere una data o un criterio di
uscita. Ogni dipendenza residua dal vecchio sistema va assegnata a un
owner e classificata: necessaria, temporanea o da eliminare.

Prima di spegnere l'origine, verificare traffico residuo, job
schedulati, webhook, client sconosciuti, backup, requisiti di audit e
retention. Conservare solo ciò che è richiesto e con accessi limitati;
spegnere credenziali e endpoint non più necessari. La riduzione della
superficie d'attacco e del costo è parte del risultato della migrazione.

\section{5.13 Approfondire la mappa dai principali
ecosistemi}\label{approfondire-la-mappa-dai-principali-ecosistemi}

I confronti tra ecosistemi diventano utili quando partono da
un'applicazione concreta. Questa sezione raccoglie differenze ricorrenti
che una migrazione dovrebbe provare esplicitamente. Non sono regole
assolute: framework e versioni cambiano, e un progetto può adottare
convenzioni diverse da quelle più comuni.

\subsection{5.13.1 Da PHP e framework PHP}\label{da-php-e-framework-php}

In una classica applicazione PHP, il modello di esecuzione può essere
per richiesta: un web server inoltra la richiesta al runtime, che carica
codice e bootstrap. Framework come Laravel o Symfony introducono
container, middleware, routing, ORM, sessioni e console command. In .NET
il processo web rimane tipicamente attivo e gestisce molte richieste
durante la propria vita.

Questa differenza impone di cercare stato globale mutabile e risorse
inizializzate una sola volta. Un singleton .NET non viene ricreato per
ogni richiesta; non deve trattenere l'utente corrente o un
\texttt{DbContext}. L'applicazione deve essere thread-safe e non
affidarsi al teardown implicito che il processo PHP poteva offrire.

Mappare:

\begin{itemize}
\tightlist
\item
  route e controller PHP verso endpoint/controller ASP.NET Core,
  conservando metodi, status e formati;
\item
  service container e service providers verso registrazioni DI con
  lifetime appropriato;
\item
  middleware e kernel events verso pipeline ASP.NET Core, distinguendo
  codice globale da logica endpoint;
\item
  Eloquent o Doctrine verso EF Core o SQL esplicito, verificando lazy
  loading, transazioni e naming;
\item
  queue worker e scheduler verso broker, worker o orchestratore, non
  necessariamente \texttt{BackgroundService};
\item
  sessioni e file upload verso storage distribuito e limiti HTTP
  espliciti.
\end{itemize}

Un aspetto spesso trascurato è il tipo dinamico dei dati. PHP può
convertire implicitamente stringhe numeriche, valori booleani o date;
C\# richiede parsing e conversioni più visibili. Non replicare
coercizioni ambigue senza capire quali client ne dipendano. Per
importare valori storici, definire cultura, decimal separator, encoding
e regole per stringhe vuote.

\subsection{5.13.2 Da Python e framework
ASGI/WSGI}\label{da-python-e-framework-asgiwsgi}

In Python distinguere il modello WSGI dal modello ASGI e capire se
l'applicazione è davvero asincrona. Un endpoint \texttt{async} che
chiama una libreria sincrona può bloccare l'event loop. In .NET,
chiamare un metodo sincrono di I/O in una route non viene reso asincrono
dalla firma del delegate: occorre usare API asincrone della libreria o
isolare il lavoro secondo un piano di capacità.

Django, Flask e FastAPI offrono convenzioni diverse per middleware,
dependency override, serializzazione, migrazioni e background task.
FastAPI genera OpenAPI da tipi e dichiarazioni, ma la specifica
esistente va verificata contro le response effettive. Non assumere che
il target generi gli stessi schema, descrizioni o default.

Per la persistenza, confrontare SQLAlchemy, Django ORM e l'eventuale SQL
diretto con il modello EF Core. Verificare transazioni implicite,
autoflush, session lifecycle, timezone e query lazy. Un Python
\texttt{None} serializzato come \texttt{null} può diventare una
proprietà assente se il serializer .NET è configurato diversamente.

Python consente spesso di distribuire sorgenti e dipendenze in un
ambiente virtuale. In .NET l'artefatto publish, il runtime e il
container sono parte della consegna. La pipeline dovrebbe ricreare
l'output da checkout pulito, non riutilizzare un \texttt{venv} o un
output locale. Le dipendenze native Python hanno analoghi problemi di
piattaforma nelle librerie .NET con componenti nativi: provare
l'immagine target, non soltanto il test locale.

\subsection{5.13.3 Da JavaScript/Node.js}\label{da-javascriptnode.js}

Node.js usa un event loop e un modello di I/O non bloccante. Il codice
CPU-bound può bloccare l'event loop; in ASP.NET Core un lavoro CPU-bound
può occupare thread del pool e saturare l'applicazione. In entrambi i
casi servono limiti, profiling e un eventuale worker separato. Non
trasferire un \texttt{setTimeout} o una promise non attesa come se
fossero una coda affidabile.

Express compone middleware con convenzioni proprie, mentre NestJS
introduce moduli, guard, pipe e interceptor. Mappare ciascuna
responsabilità, senza trasformare ogni middleware in un servizio DI o
ogni provider Nest in singleton .NET. Controllare l'ordine di
middleware, la gestione degli errori asincroni e l'effetto delle route
non corrispondenti.

TypeScript fornisce controlli statici durante la build ma i dati HTTP
restano dinamici a runtime. C\# offre un sistema statico integrato nel
compilatore, tuttavia il binding JSON non sostituisce validazione. Non
rimuovere i controlli runtime perché il DTO ha proprietà non-nullable:
client non conformi possono inviare body incompleti.

Controllare differenze di serializzazione: \texttt{undefined} in
JavaScript può essere omesso, \texttt{null} è distinto; in JSON non
esiste \texttt{undefined}. Le opzioni di naming camelCase, enum come
stringhe o numeri, date e proprietà ignorate devono essere definite e
testate. Verificare inoltre gli stream: bufferizzare tutto il body può
aumentare memoria e latenza rispetto a una pipeline che processava dati
incrementali.

\subsection{5.13.4 Da Java e Spring Boot}\label{da-java-e-spring-boot}

Java e .NET condividono alcuni concetti di runtime gestito, garbage
collection, compilazione e applicazioni server long-running, ma il
modello di librerie e framework resta specifico. In Spring, confrontare
bean scope, filtri servlet, interceptor, controller advice,
configurazione e transazioni con i componenti ASP.NET Core
corrispondenti.

Un bean \texttt{singleton} Spring e un singleton DI .NET hanno entrambi
una durata ampia, ma il comportamento effettivo dipende da proxy, scope
e lifecycle del framework. Controllare le dipendenze scoped e la
sicurezza concorrente. Non assumere che un \texttt{@Transactional} si
traduca in una transazione che attraversa chiamate HTTP o eventi.

Confrontare JPA/Hibernate con EF Core su tracking, flush, lazy loading,
cascade, lock ottimistico e query generation. Le due tecnologie possono
produrre SQL diverso e differire nei limiti del provider. Provare query
e migrazioni sullo stesso motore database; una conversione di annotation
non è un piano dati.

Confrontare Maven/Gradle con NuGet e i file \texttt{.csproj} riguardo a
versioni, dipendenze transitive, plugin e packaging. Verificare profili
di build, variabili ambientali e configurazione del server. I test su
Spring Test o container di test vanno sostituiti da test equivalenti sul
target e sul provider effettivo, non semplicemente rimossi.

\subsection{5.13.5 Tabella di mapping delle
responsabilità}\label{tabella-di-mapping-delle-responsabilituxe0}

Durante l'assessment costruire una tabella concreta. Esempio:

{\def\LTcaptype{none} % do not increment counter
\begin{longtable}[]{@{}
  >{\raggedright\arraybackslash}p{(\linewidth - 6\tabcolsep) * \real{0.2500}}
  >{\raggedright\arraybackslash}p{(\linewidth - 6\tabcolsep) * \real{0.2500}}
  >{\raggedright\arraybackslash}p{(\linewidth - 6\tabcolsep) * \real{0.2500}}
  >{\raggedright\arraybackslash}p{(\linewidth - 6\tabcolsep) * \real{0.2500}}@{}}
\toprule\noalign{}
\begin{minipage}[b]{\linewidth}\raggedright
Sistema sorgente
\end{minipage} & \begin{minipage}[b]{\linewidth}\raggedright
Comportamento osservato
\end{minipage} & \begin{minipage}[b]{\linewidth}\raggedright
Candidato .NET
\end{minipage} & \begin{minipage}[b]{\linewidth}\raggedright
Prova di migrazione
\end{minipage} \\
\midrule\noalign{}
\endhead
\bottomrule\noalign{}
\endlastfoot
Middleware Laravel & Aggiunge tenant da un claim e blocca utenti sospesi
& Middleware più policy o servizio endpoint & Test anonimo, utente
sospeso e tenant valido \\
Dependency FastAPI & Fornisce repository per richiesta & Servizio scoped
e repository/query service & Due richieste non condividono stato \\
Middleware Express & Trasforma errori asincroni in JSON & Exception
handler centrale & Eccezioni note e inattese con status coerente \\
Filtro Spring & Verifica permesso di scrittura & Policy authorization &
Token scaduto, claim mancante e claim valido \\
Job Celery & Import idempotente con retry & Worker e broker/job
scheduler & Crash durante batch e replay senza duplicati \\
Migrazione Doctrine & Aggiunge campo con backfill & Script controllato
expand--migrate--contract & Rollout con vecchia versione ancora
attiva \\
\end{longtable}
}

La colonna ``candidato .NET'' è un'ipotesi finché la prova non la
conferma. Annotare anche alternative e decisione. Se il comportamento
sorgente non è determinabile, creare prima un test o una sessione di
discovery.

\subsection{5.13.6 Scheda di decisione per un
concetto}\label{scheda-di-decisione-per-un-concetto}

Una scheda riutilizzabile può avere questi campi:

\begin{Shaded}
\begin{Highlighting}[]
\NormalTok{Concetto sorgente:}
\NormalTok{Versione e configurazione:}
\NormalTok{Responsabilità osservata:}
\NormalTok{Input e output:}
\NormalTok{Lifecycle e scope:}
\NormalTok{Dipendenze e side effect:}
\NormalTok{Candidato ASP.NET Core:}
\NormalTok{Differenze semantiche:}
\NormalTok{Test che dimostra il mapping:}
\NormalTok{Decisione e motivazione:}
\NormalTok{Owner e data di revisione:}
\end{Highlighting}
\end{Shaded}

La scheda impedisce due errori opposti: tradurre tutto meccanicamente o
ricominciare senza preservare nulla. I dettagli si possono archiviare
vicino alla decisione tecnica e collegare ai test; il capitolo o la
documentazione di progetto deve conservare almeno la motivazione e le
incompatibilità ancora aperte.

\subsection{5.13.7 Quando mantenere un componente nel sistema
originario}\label{quando-mantenere-un-componente-nel-sistema-originario}

Non tutto deve essere migrato immediatamente. Un motore di report con
dipendenze non portabili, un adapter verso hardware o un job stabile
possono restare temporaneamente dietro un contratto di servizio. La
decisione va confrontata con il costo di mantenere runtime, competenze,
security patch e pipeline separate.

Un componente mantenuto deve avere:

\begin{itemize}
\tightlist
\item
  owner, responsabile degli aggiornamenti e processo di incident
  response;
\item
  interfaccia di rete o messaggio documentata, con autenticazione e
  timeout;
\item
  monitoraggio e piano di continuità;
\item
  dipendenze e dati che ne definiscono la vita residua;
\item
  criterio e data di revisione della decisione.
\end{itemize}

La coesistenza è una strategia, non un risultato finale implicito. Se la
migrazione completa non è conveniente, dichiarare il confine come
architettura permanente e finanziarne la manutenzione; se è transitorio,
misurarne la riduzione.

\section{5.14 Prove di migrazione dei dati: runbook e
controlli}\label{prove-di-migrazione-dei-dati-runbook-e-controlli}

Lo spostamento dei dati è un processo controllato con input, checkpoint
e report. Il runbook seguente aiuta a trasformare uno script una tantum
in un'attività ripetibile, verificabile e recuperabile.

\subsection{5.14.1 Preparare il runbook}\label{preparare-il-runbook}

Prima di eseguire una copia reale, documentare:

\begin{enumerate}
\def\labelenumi{\arabic{enumi}.}
\tightlist
\item
  database sorgente e target, versioni, region e responsabili;
\item
  oggetti e intervalli inclusi, esclusioni e regole di trasformazione;
\item
  permessi minimi delle identità usate e scadenza delle credenziali;
\item
  volume previsto, throughput, durata stimata e impatto sulle query;
\item
  punto di consistenza, strategia delta e trattamento delle scritture
  concorrenti;
\item
  checkpoint, retry, deduplicazione e ripresa dopo interruzione;
\item
  validazioni, soglie di errore, approvatore e piano di rollback;
\item
  retention e protezione dei file temporanei e dei report.
\end{enumerate}

Il runbook non deve includere segreti in chiaro. I riferimenti alle
credenziali devono indicare il secret store e l'identità autorizzata.
Limitare accesso ai dump, cifrarli in transito e a riposo e cancellarli
secondo retention approvata.

\subsection{5.14.2 Snapshot, delta e
freeze}\label{snapshot-delta-e-freeze}

Il piano di migrazione deve definire cosa significa ``dati coerenti''.
Una fotografia estratta in più query a orari diversi può contenere righe
incompatibili tra loro se il sistema continua a scrivere. Usare snapshot
transazionali o un meccanismo CDC quando il provider lo supporta e la
finestra di consistenza è necessaria.

Se si effettua un freeze:

\begin{itemize}
\tightlist
\item
  comunicare in anticipo il blocco e identificare tutte le write path,
  incluso back office;
\item
  fermare job e consumer in ordine concordato;
\item
  attendere che le transazioni in corso si completino;
\item
  registrare watermark temporale o identificativo dell'ultimo record;
\item
  estrarre, trasformare e importare;
\item
  eseguire riconciliazione e test applicativi;
\item
  riaprire le scritture una sola volta deciso il sistema owner.
\end{itemize}

Se il freeze non è possibile, usare uno snapshot più un delta ordinato e
definire come trattare update e delete. Il timestamp
\texttt{updated\_at} da solo può non bastare se non è monotono, ha
precisione insufficiente o non registra cancellazioni. Un log
transazionale o una tabella outbox può fornire una sorgente di
cambiamenti più affidabile.

\subsection{5.14.3 Trasformazioni
versionate}\label{trasformazioni-versionate}

Ogni trasformazione deve essere identificata da una versione dello
script e avere esempi prima/dopo. Separare parsing, validazione, mapping
e scrittura per poter testare le regole senza connettersi ai database.

Esempi di trasformazioni che richiedono approvazione:

\begin{itemize}
\tightlist
\item
  stringa vuota convertita in \texttt{null};
\item
  codice paese normalizzato;
\item
  date locali interpretate in un fuso nominato;
\item
  enum legacy mappati a nuovi valori;
\item
  importi con scala differente;
\item
  record duplicati uniti secondo una regola di dominio.
\end{itemize}

Non inferire unità o timezone da un nome di campo ambiguo. Se un dato
non può essere interpretato, metterlo in quarantena con ragione e chiave
tecnica; non scartarlo silenziosamente. I report non devono esporre
valori personali se bastano conteggio, categoria di errore e
identificativo pseudonimo.

\subsection{5.14.4 Importazione
riprendibile}\label{importazione-riprendibile}

Per dataset grandi elaborare batch limitati, registrando il checkpoint
in storage durevole. La chiave di ripresa deve rappresentare un ordine
stabile, come una chiave primaria o un cursore monotono, non l'offset di
una query che cambia mentre si importano righe.

Ogni batch può seguire questo ciclo:

\begin{enumerate}
\def\labelenumi{\arabic{enumi}.}
\tightlist
\item
  leggere un intervallo consistente;
\item
  trasformare e validare in memoria limitata;
\item
  scrivere in transazione o staging;
\item
  verificare il batch;
\item
  registrare il checkpoint solo dopo il commit;
\item
  proseguire o fermarsi su soglia di errore.
\end{enumerate}

Se il processo termina dopo il commit ma prima del checkpoint, il batch
può essere ripetuto. Usare upsert con chiave stabile o una tabella di
deduplicazione; non assumere consegna exactly-once. Il checkpoint prima
del commit può invece saltare dati. Definire una strategia per ripetere
in sicurezza.

\subsection{5.14.5 Staging e cutover}\label{staging-e-cutover}

Una tabella di staging consente di importare e validare senza esporre
immediatamente i record agli utenti. Può contenere migration id,
origine, timestamp e stato di validazione. Proteggere lo staging come il
dato originale; non lasciarlo diventare un archivio permanente di copie
non governate.

Il cutover richiede una decisione sulla lettura autorevole. Possibili
percorsi: cambiare connection string, attivare un routing di tenant,
promuovere tabelle o aggiornare un alias. Documentare la cache e la
propagazione della configurazione, perché il cambio potrebbe non essere
istantaneo.

Prima dell'attivazione, congelare le scritture secondo il piano oppure
completare il delta, validare i criteri e verificare accesso e
performance. Dopo l'attivazione, monitorare errori e mismatch per un
periodo concordato; mantenere l'origine in sola lettura finché il
rollback dati è ancora possibile.

\subsection{5.14.6 Report di
riconciliazione}\label{report-di-riconciliazione}

Un report utile include:

{\def\LTcaptype{none} % do not increment counter
\begin{longtable}[]{@{}
  >{\raggedright\arraybackslash}p{(\linewidth - 6\tabcolsep) * \real{0.2143}}
  >{\raggedleft\arraybackslash}p{(\linewidth - 6\tabcolsep) * \real{0.2857}}
  >{\raggedleft\arraybackslash}p{(\linewidth - 6\tabcolsep) * \real{0.2857}}
  >{\raggedright\arraybackslash}p{(\linewidth - 6\tabcolsep) * \real{0.2143}}@{}}
\toprule\noalign{}
\begin{minipage}[b]{\linewidth}\raggedright
Verifica
\end{minipage} & \begin{minipage}[b]{\linewidth}\raggedleft
Origine
\end{minipage} & \begin{minipage}[b]{\linewidth}\raggedleft
Target
\end{minipage} & \begin{minipage}[b]{\linewidth}\raggedright
Esito
\end{minipage} \\
\midrule\noalign{}
\endhead
\bottomrule\noalign{}
\endlastfoot
Righe lette, scritte, scartate & conteggio & conteggio &
riconciliato? \\
Record per stato e tenant & aggregato & aggregato & scostamento \\
Chiavi orfane & conteggio & conteggio & soglia \\
Hash di campioni canonici & hash & hash & corrispondenza \\
Valori non trasformabili & conteggio & quarantena & owner \\
\end{longtable}
}

I checksum devono essere calcolati su una rappresentazione canonica che
definisce ordinamento, encoding e null. Un hash dell'intero file può
rilevare differenze, ma non spiega quale record o trasformazione le
abbia causate. Combinare aggregati, campioni e conteggi per localizzare
rapidamente il problema.

Il report include migration id, commit dello script, configurazione non
segreta, orari, watermark, durata e identità esecutrice. Conservare il
report secondo policy e rendere riproducibili le query di verifica.

\subsection{5.14.7 Criteri di stop}\label{criteri-di-stop}

Interrompere l'importazione o il cutover se:

\begin{itemize}
\tightlist
\item
  gli errori superano la soglia concordata o coinvolgono dati critici;
\item
  il delta non ha un watermark affidabile;
\item
  la replica o il database supera il budget di carico;
\item
  gli aggregati divergono senza una regola approvata;
\item
  non è chiaro quale sistema possiede le scritture;
\item
  il rollback richiesto non è più praticabile.
\end{itemize}

Fermarsi è un esito valido. Il runbook deve indicare come lasciare il
sistema in uno stato consistente e chi decide se riprendere,
ripristinare o applicare una correzione.

\section{5.15 Diagnosi operativa durante la
migrazione}\label{diagnosi-operativa-durante-la-migrazione}

Gli incidenti di migrazione spesso emergono solo quando il target
incontra rete, proxy, credenziali e carichi reali. Una diagnosi efficace
parte dai segnali e dal percorso della richiesta, non dalla modifica
casuale di configurazioni.

\subsection{5.15.1 La readiness fallisce, ma il processo è
attivo}\label{la-readiness-fallisce-ma-il-processo-uxe8-attivo}

Separare liveness da readiness e leggere il dettaglio del controllo che
fallisce. Verificare che hostname, certificati, DNS, accessi database e
timeout siano quelli dell'ambiente target. Un health check che dipende
da tutte le API esterne può rendere indisponibile l'intero servizio per
una dipendenza non critica.

Confrontare configurazione attiva e sorgente attesa senza stampare
segreti. Registrare quali provider di configurazione sono abilitati e
validare opzioni all'avvio. Se il servizio si avvia con un endpoint
esterno vuoto o un valore di default, correggere la validazione invece
di aggiungere retry indefiniti.

\subsection{5.15.2 Errori 401 o 403 dopo il
cutover}\label{errori-401-o-403-dopo-il-cutover}

401 indica in genere un problema di autenticazione; 403 indica identità
riconosciuta ma non autorizzata, anche se il contratto può variare.
Verificare emittente (\texttt{iss}), destinatario (\texttt{aud}),
scadenze, chiavi di firma, scarto orario e claim ricevuti. Confrontare
il token in modo sicuro senza copiarlo nei ticket o nei log.

Controllare mapping di claim e differenze di casing o nomi tra provider.
Un framework sorgente può convertire automaticamente gruppi in ruoli,
mentre il target richiede una policy esplicita. Testare un utente
rappresentativo per ogni profilo e distinguere permesso mancante da
route non autenticata. Non disabilitare authorization per ``sbloccare il
deploy''.

\subsection{5.15.3 502, 503 e timeout
intermittenti}\label{e-timeout-intermittenti}

Seguire la richiesta attraverso client, gateway, ingress, Kestrel e
dipendenze. Confrontare timeout configurati a ogni livello e verificare
se un proxy chiude la connessione prima dell'applicazione. Correlare
trace id e timestamp tra sistemi; senza sincronizzazione oraria i log
possono sembrare contraddittori.

Osservare pool connessioni, thread pool, code, DNS e saturazione
database. Un retry su un timeout può moltiplicare il carico e peggiorare
l'incidente. Ridurre temporaneamente il traffico o disattivare la
capability può essere più sicuro che aumentare tutti i timeout.

\subsection{5.15.4 Differenze di
serializzazione}\label{differenze-di-serializzazione}

Se i test passano ma un client rifiuta la response, confrontare byte e
schema: \texttt{Content-Type}, casing, proprietà omesse, enum, formato
data, precisione numerica e encoding. Verificare opzioni globali del
serializer e attributi sul DTO. Un aggiornamento del serializer o un
modello con nullability diversa può cambiare il JSON senza modificare la
route.

Riprodurre con fixture minime e un test contrattuale. Non patchare la
risposta in middleware per correggere un singolo DTO: la regola deve
vivere nel mapping o nella configurazione condivisa documentata.

\subsection{5.15.5 Mismatch di database dopo la
migrazione}\label{mismatch-di-database-dopo-la-migrazione}

Prima di correggere i record, verificare query, filtri, timezone,
collations, precisione e ordine. Calcolare differenze per categoria e
chiave, usando dati minimizzati. Una differenza può provenire da una
query N+1 o da una proprietà null, non dalla copia dati.

Confrontare SQL generato e piani d'esecuzione. Verificare se il vecchio
ORM applicava soft-delete, tenant filter o conversioni implicite. Non
riparare manualmente il database live prima di sapere se il processo di
importazione può riprodurre il problema.

\subsection{5.15.6 Runbook di incidente}\label{runbook-di-incidente}

Per un errore in rollout:

\begin{enumerate}
\def\labelenumi{\arabic{enumi}.}
\tightlist
\item
  identificare release, percentuale traffico e capability coinvolta;
\item
  verificare alert e impatto utente, distinguendo sintomo e causa;
\item
  applicare la soglia predefinita: fermare incremento o attivare
  rollback;
\item
  proteggere dati e fermare eventuali writer duplicati;
\item
  salvare log e metriche necessari, redigendo segreti e dati personali;
\item
  assegnare un responsabile tecnico e un comunicatore;
\item
  dopo la stabilizzazione, documentare causa, rilevazione e prevenzione.
\end{enumerate}

Non fare debug direttamente in produzione con logging verboso permanente
o query non controllate. Ogni modifica temporanea ha owner, scadenza e
piano di rimozione.

\section{5.16 Verifica finale: criteri di
successo}\label{verifica-finale-criteri-di-successo}

Considerare la migrazione conclusa solo quando il target possiede
chiaramente le capability previste e il sistema sorgente non mantiene
dipendenze operative nascoste. Account, job, backup, credenziali e
supporto devono avere un owner nel nuovo assetto o una data di
dismissione.

Misurare il successo con criteri concordati: flussi critici conformi,
error rate e latenza entro budget, dati riconciliati, sicurezza
approvata, rollback provato, operatori formati e costi del vecchio
sistema rimossi o motivati. Conservare le decisioni e i report in modo
accessibile al team che sarà on-call.

\subsection{5.16.1 Fonti ufficiali e
aggiornamento}\label{fonti-ufficiali-e-aggiornamento}

Consultare documentazione ufficiale e controllare la versione target
prima di applicare un esempio. Le pagine evolvono; per un progetto reale
annotare la data di verifica e la versione .NET:

\begin{itemize}
\tightlist
\item
  \href{https://learn.microsoft.com/dotnet/core/porting/}{.NET ---
  Porting to .NET}
\item
  \href{https://dotnet.microsoft.com/platform/support/policy/dotnet-core}{.NET
  --- Panoramica del ciclo di vita}
\item
  \href{https://learn.microsoft.com/aspnet/core/introduction-to-aspnet-core}{ASP.NET
  Core --- Panoramica}
\item
  \href{https://learn.microsoft.com/aspnet/core/fundamentals/middleware/}{ASP.NET
  Core --- Middleware}
\item
  \href{https://learn.microsoft.com/dotnet/core/extensions/dependency-injection}{ASP.NET
  Core --- Dependency Injection}
\item
  \href{https://learn.microsoft.com/aspnet/core/fundamentals/configuration/}{ASP.NET
  Core --- Configurazione}
\item
  \href{https://learn.microsoft.com/aspnet/core/security/authentication/}{ASP.NET
  Core --- Autenticazione}
\item
  \href{https://learn.microsoft.com/aspnet/core/test/integration-tests}{ASP.NET
  Core --- Test d'integrazione}
\item
  \href{https://learn.microsoft.com/ef/core/managing-schemas/migrations/}{Entity
  Framework Core --- Migrazioni}
\item
  \href{https://learn.microsoft.com/ef/core/performance/}{Entity
  Framework Core --- Performance}
\item
  \href{https://learn.microsoft.com/dotnet/core/docker/}{.NET ---
  Docker}
\item
  \href{https://learn.microsoft.com/dotnet/architecture/}{.NET ---
  Architettura delle applicazioni}
\end{itemize}

\section{5.12 Caso pratico: migrare un servizio
catalogo}\label{caso-pratico-migrare-un-servizio-catalogo}

Per collegare le decisioni, immaginiamo di migrare un servizio catalogo
esistente, scritto in Python con FastAPI, verso ASP.NET Core. Il
servizio pubblica prodotti, legge un database PostgreSQL condiviso con
un back office e viene consumato da un sito e da un'app mobile.
L'esempio non è un tutorial per ricostruire un catalogo completo: mostra
come delimitare il lavoro, mantenere il contratto e produrre evidenze
prima di spostare il traffico.

\subsection{5.12.1 Scenario e vincoli}\label{scenario-e-vincoli}

Il sistema attuale espone:

\begin{itemize}
\tightlist
\item
  \texttt{GET\ /api/products?category=...\&page=...} per elencare
  prodotti attivi;
\item
  \texttt{GET\ /api/products/\{id\}} per leggere un prodotto;
\item
  \texttt{POST\ /api/products} per creare un prodotto, riservato agli
  editor;
\item
  un job notturno che importa disponibilità da un partner.
\end{itemize}

Il database contiene milioni di righe storiche, mentre l'endpoint
pubblico restituisce solo prodotti attivi. Il back office continua a
scrivere nello schema durante la migrazione. L'app mobile non può essere
aggiornata rapidamente. La disponibilità richiesta non consente un fermo
lungo e l'autenticazione usa un identity provider aziendale.

L'assessment mostra che le letture pubbliche sono il primo candidato:
sono frequenti, hanno contratto documentabile e possono rimanere
read-only sul database esistente. La scrittura di prodotti resta sul
servizio sorgente nella prima fase; non si duplicano gli effetti di
\texttt{POST}. Il job notturno viene migrato separatamente perché ha un
lifecycle e un profilo di retry diversi.

\subsection{5.12.2 Baseline e test di
caratterizzazione}\label{baseline-e-test-di-caratterizzazione}

Prima di scrivere il target si raccolgono:

\begin{itemize}
\tightlist
\item
  esempi di response per prodotto attivo, assente e non visibile;
\item
  comportamento dei filtri e paginazione;
\item
  status e schema per query non valide;
\item
  requisiti di autorizzazione del \texttt{POST};
\item
  p95/p99, richieste al secondo, error rate e carico sul database;
\item
  piano query corrente e indici usati.
\end{itemize}

Le fixture contengono dati sintetici con casi di null, descrizione
lunga, prezzo al limite e categoria inesistente. Per i casi di
produzione si usano payload anonimizzati, senza identificativi o
informazioni personali. I test memorizzano anche ordinamento e proprietà
JSON richieste dai client.

Si concorda che l'implementazione nuova deve restituire la stessa
semantica dei filtri e i medesimi codici, salvo variazioni approvate. Il
nuovo sistema deve mantenere il budget di latenza del servizio e non
superare il limite di connessioni assegnato al database.

\subsection{5.12.3 Creazione del progetto}\label{creazione-del-progetto}

La solution può iniziare con API e test, senza creare tutti i progetti
architetturali possibili:

\begin{Shaded}
\begin{Highlighting}[]
\NormalTok{Catalogo.sln}
\NormalTok{src/}
\NormalTok{  Catalogo.Api/}
\NormalTok{  Catalogo.Infrastructure/}
\NormalTok{tests/}
\NormalTok{  Catalogo.Api.Tests/}
\end{Highlighting}
\end{Shaded}

Il progetto API configura hosting, autenticazione, logging, endpoint e
composizione dei servizi. Infrastructure contiene \texttt{DbContext},
mapping e query SQL/EF Core. Se le regole del dominio diventano
complesse, si può introdurre un progetto Domain con una decisione
esplicita, non come prerequisito per un endpoint di lettura.

La foundation si considera pronta quando un commit pulito passa
\texttt{dotnet\ restore}, \texttt{dotnet\ build}, \texttt{dotnet\ test},
produce l'immagine, gira nel namespace di test e legge configurazione
esterna. Un health check minimale può essere registrato e mappato così:

\begin{Shaded}
\begin{Highlighting}[]
\DataTypeTok{var}\NormalTok{ builder }\OperatorTok{=}\NormalTok{ WebApplication}\OperatorTok{.}\FunctionTok{CreateBuilder}\OperatorTok{(}\NormalTok{args}\OperatorTok{);}
\NormalTok{builder}\OperatorTok{.}\FunctionTok{Services}\OperatorTok{.}\FunctionTok{AddHealthChecks}\OperatorTok{();}

\DataTypeTok{var}\NormalTok{ app }\OperatorTok{=}\NormalTok{ builder}\OperatorTok{.}\FunctionTok{Build}\OperatorTok{();}
\NormalTok{app}\OperatorTok{.}\FunctionTok{MapHealthChecks}\OperatorTok{(}\StringTok{"/health"}\OperatorTok{);}
\NormalTok{app}\OperatorTok{.}\FunctionTok{Run}\OperatorTok{();}
\end{Highlighting}
\end{Shaded}

Questo endpoint verifica che l'applicazione risponda; controlli di
readiness sulle dipendenze vanno aggiunti solo se richiesti. La pipeline
non usa una connection string reale nei log.

\subsection{5.12.4 Contratto e query di
lettura}\label{contratto-e-query-di-lettura}

Il DTO è separato dal modello persistito:

\begin{Shaded}
\begin{Highlighting}[]
\KeywordTok{public} \KeywordTok{sealed}\NormalTok{ record }\FunctionTok{ProductListItem}\OperatorTok{(}
    \DataTypeTok{int}\NormalTok{ Id}\OperatorTok{,}
    \DataTypeTok{string}\NormalTok{ Name}\OperatorTok{,}
    \DataTypeTok{decimal}\NormalTok{ Price}\OperatorTok{,}
    \DataTypeTok{string}\NormalTok{ Currency}\OperatorTok{);}

\KeywordTok{public} \KeywordTok{sealed}\NormalTok{ record }\FunctionTok{ProductPage}\OperatorTok{(}
\NormalTok{    IReadOnlyList}\OperatorTok{\textless{}}\NormalTok{ProductListItem}\OperatorTok{\textgreater{}}\NormalTok{ Items}\OperatorTok{,}
    \DataTypeTok{int}\NormalTok{ Page}\OperatorTok{,}
    \DataTypeTok{int}\NormalTok{ PageSize}\OperatorTok{,}
    \DataTypeTok{int}\NormalTok{ TotalCount}\OperatorTok{);}
\end{Highlighting}
\end{Shaded}

La paginazione deve essere limitata lato server; non accettare page size
illimitata. Per dataset in crescita, una paginazione a cursore può
essere più stabile ed efficiente di un offset elevato, ma cambiare il
contratto va coordinato con i client. Nella prima fase si mantiene
l'interfaccia esistente e si usa il piano database per verificare il
costo.

Un service applicativo riceve filtri già validati e restituisce un
risultato; l'endpoint mappa il risultato a HTTP:

\begin{Shaded}
\begin{Highlighting}[]
\NormalTok{app}\OperatorTok{.}\FunctionTok{MapGet}\OperatorTok{(}\StringTok{"/api/products"}\OperatorTok{,} \FunctionTok{async} \OperatorTok{(}
    \DataTypeTok{string}\OperatorTok{?}\NormalTok{ category}\OperatorTok{,}
    \DataTypeTok{int}\NormalTok{ page}\OperatorTok{,}
    \DataTypeTok{int}\NormalTok{ pageSize}\OperatorTok{,}
\NormalTok{    IProductQueries queries}\OperatorTok{,}
\NormalTok{    CancellationToken ct}\OperatorTok{)} \OperatorTok{=\textgreater{}}
\OperatorTok{\{}
    \KeywordTok{if} \OperatorTok{(}\NormalTok{page }\OperatorTok{\textless{}} \DecValTok{1} \OperatorTok{||}\NormalTok{ pageSize }\KeywordTok{is} \OperatorTok{\textless{}} \DecValTok{1}\NormalTok{ or }\OperatorTok{\textgreater{}} \DecValTok{100}\OperatorTok{)}
        \KeywordTok{return}\NormalTok{ Results}\OperatorTok{.}\FunctionTok{BadRequest}\OperatorTok{();}

    \DataTypeTok{var}\NormalTok{ result }\OperatorTok{=}\NormalTok{ await queries}\OperatorTok{.}\FunctionTok{ListAsync}\OperatorTok{(}
\NormalTok{        category}\OperatorTok{,}\NormalTok{ page}\OperatorTok{,}\NormalTok{ pageSize}\OperatorTok{,}\NormalTok{ ct}\OperatorTok{);}
    \KeywordTok{return}\NormalTok{ Results}\OperatorTok{.}\FunctionTok{Ok}\OperatorTok{(}\NormalTok{result}\OperatorTok{);}
\OperatorTok{\});}
\end{Highlighting}
\end{Shaded}

Il frammento illustra il flusso e non definisce da solo tutte le regole:
i default per i parametri, il formato errore, i limiti e l'ordinamento
vanno allineati al contratto sorgente. In codice di produzione la
validazione restituisce una risposta coerente e documentata; l'endpoint
non deve trasformare un errore database in lista vuota.

\subsection{5.12.5 Mantenere il database condiviso in
sicurezza}\label{mantenere-il-database-condiviso-in-sicurezza}

Il nuovo servizio apre un accesso read-only se il provider e i permessi
lo consentono. Questo riduce il rischio di alterare dati durante la
prima fase. Le query sono proiettate sui campi necessari, limitano
risultati e usano parametri. Si controllano indici e piani; l'utente DB
non ha permessi DDL o scrittura non necessari.

Prima del deploy si verificano schema e tipi con i proprietari del back
office. Se il vecchio sistema cambia una colonna, il target deve fallire
con log diagnostico e alert, non continuare con un risultato parziale
apparentemente valido. Un contratto dati o una vista stabile può isolare
il nuovo servizio dalle tabelle interne, se il team database può
mantenerli.

Il modello EF Core mappa solo la parte necessaria dello schema, senza
avviare migrazioni automatiche. Il target non è proprietario del
database in questa fase; non può quindi presumere di aggiornare tabelle
condivise. Le migrazioni future richiedono un piano separato con tutti i
writer coinvolti.

\subsection{5.12.6 Test differenziale e osservazione in
shadow}\label{test-differenziale-e-osservazione-in-shadow}

La suite d'integrazione invia richieste al sistema sorgente e al target
usando un database di test fissato. Le risposte sono confrontate su
status, schema e valori. Timestamp e campi variabili sono normalizzati
solo dove il contratto li dichiara non deterministici.

In ambiente di preproduzione il traffico sintetico esercita route e
limiti. In produzione, se privacy, policy e architettura lo permettono,
il gateway può duplicare \textbf{solo richieste di lettura} verso il
target in modalità shadow. Il target esegue il calcolo ma la sua
risposta non viene restituita all'utente. Si confrontano risultati in
forma redatta e si misurano query e latenza.

Non duplicare il \texttt{POST} o il job di importazione: un secondo
invio potrebbe creare un prodotto, inviare notifiche o alterare stock.
Shadow traffic non è una licenza a replicare dati personali; minimizzare
payload, applicare retention breve e ottenere le approvazioni
necessarie.

Un report giornaliero classifica differenze per categoria: status,
proprietà mancanti, ordinamento, valori numerici, null e timeout. Ogni
differenza ha owner e decisione: bug nel nuovo codice, bug sorgente da
correggere, differenza intenzionale o dato invalido. Il traffico non
viene spostato finché le differenze critiche non sono comprese.

\subsection{5.12.7 Canary e cutover}\label{canary-e-cutover}

Con test differenziali soddisfacenti, il gateway indirizza al target una
piccola percentuale di richieste di lettura, mantenendo identificabile
la coorte. Si controllano:

\begin{itemize}
\tightlist
\item
  error rate e status per route e versione;
\item
  p95/p99 e timeout;
\item
  pool connessioni, CPU, memoria e query lente;
\item
  differenze di contenuto e errori di deserializzazione client;
\item
  log di autenticazione e richieste senza tenant.
\end{itemize}

Si amplia gradualmente la percentuale solo se le metriche restano entro
budget per una finestra definita. Il criterio non è ``nessuna
segnalazione'': un client mobile può essere usato poco e segnalare
tardi. Gli alert e i log devono offrire copertura tecnica indipendente
dal feedback degli utenti.

Per rollback, il gateway riporta la route al servizio Python. Poiché il
target è read-only e usa lo schema invariato, il rollback non richiede
sincronizzare nuove scritture. Si prova prima in preproduzione,
verificando cache, DNS, warm-up e propagazione della configurazione. Il
piano include un referente che può arrestare il rollout.

\subsection{5.12.8 Migrare la scrittura}\label{migrare-la-scrittura}

Una volta stabile la lettura, la creazione prodotto diventa una capacità
separata. Il team decide se il nuovo servizio assume la scrittura oppure
se il vecchio sistema resta owner temporaneamente. Non basta duplicare
le scritture: si deve evitare che back office e API sovrascrivano
modifiche in modo inconsistente.

Se il target assume ownership, il rollout della scrittura avviene solo
dopo aver verificato:

\begin{enumerate}
\def\labelenumi{\arabic{enumi}.}
\tightlist
\item
  policy di autorizzazione e mapping claims;
\item
  validazione e invarianti condivise con il back office;
\item
  unicità e concorrenza nel database;
\item
  audit e notifiche prodotte una sola volta;
\item
  idempotenza delle richieste ritentate;
\item
  schema compatibile con la versione precedente durante il rollback.
\end{enumerate}

Una chiave di idempotenza può aiutare per richieste ripetibili, ma deve
avere storage e scadenza definiti. L'outbox può garantire che un evento
venga registrato nella stessa transazione della creazione, con consegna
asincrona e consumer deduplicante. Non rende transazionale il sistema
esterno.

\subsection{5.12.9 Migrare il job
notturno}\label{migrare-il-job-notturno}

Il job è una capacità distinta dall'API. Si verificano frequenza,
finestra, volume, checkpoint, dati parziali e comportamento al retry. Un
\texttt{BackgroundService} può ospitare un ciclo semplice, ma se il job
deve sopravvivere a riavvii, essere schedulato da più istanze o
riprendere da un checkpoint, si usa un orchestratore o una coda
durevole.

Durante la prova, il target importa in un ambiente isolato e confronta
conteggi e aggregati con il job esistente. Non si eseguono entrambi in
scrittura sullo stesso dataset senza una strategia d'idempotenza. Il
cutover del job ha un lock o un singolo scheduler owner, un run id,
report e procedura di replay.

\subsection{5.12.10 Risultato e lezioni}\label{risultato-e-lezioni}

Il servizio di lettura può essere considerato migrato quando il target
ha sostenuto il traffico previsto, rispetta contratto e budget, è
osservabile, ha un rollback provato e il sistema sorgente non riceve più
quelle route. Il capitolo di progetto non si chiude però finché
scrittura, job, integrazioni residue e dismissione non hanno owner e
date.

Le decisioni hanno mantenuto il perimetro piccolo, ma non hanno ignorato
il sistema completo. Si è iniziato con letture reversibili; dati
condivisi e scritture sono rimasti sotto un owner; test e telemetria
hanno preceduto il cutover; infine ogni capability ha avuto il suo
piano. In un caso PHP o Java cambierebbero framework, tool e dettagli di
persistenza, non il metodo di mappatura, caratterizzazione e rollout.

\subsection{5.12.11 Adattare il caso ad altri
ecosistemi}\label{adattare-il-caso-ad-altri-ecosistemi}

\begin{itemize}
\tightlist
\item
  \textbf{PHP/Laravel o Symfony:} inventariare middleware, service
  container, queue worker, session driver e convenzioni
  Eloquent/Doctrine. Verificare differenze di null, timezone e gestione
  dei job.
\item
  \textbf{Python/Django o FastAPI:} distinguere ASGI/WSGI, dependency
  overrides, task asincroni, ORM e migration framework. Non assumere che
  codice \texttt{async} sia non bloccante se una libreria sottostante è
  sincrona.
\item
  \textbf{Node.js/Express o NestJS:} mappare middleware e
  pipe/interceptor, controllare event loop, stream e dipendenze native
  dei pacchetti.
\item
  \textbf{Java/Spring Boot:} confrontare lifecycle bean, scope request,
  filtri, transazioni e thread pool. Verificare semantica di ORM, lazy
  loading e gestione delle eccezioni.
\end{itemize}

Per ciascuno, costruire una scheda per concetto: comportamento sorgente,
candidato ASP.NET Core, differenza osservabile, test di prova, impatto e
decisione. Un mapping documentato riduce il rischio che una differenza
venga scoperta solo da un utente in produzione.

\section{5.17 Checklist, glossario e
anti-pattern}\label{checklist-glossario-e-anti-pattern}

Una migrazione efficace è una sequenza di decisioni verificabili. Le
checklist che seguono non sostituiscono il giudizio del team: aiutano a
non dimenticare dipendenze e criteri prima di una fase irreversibile.

\subsection{5.17.1 Checklist di
assessment}\label{checklist-di-assessment}

\begin{itemize}
\tightlist
\item[$\square$]
  Sono stati identificati owner funzionali e tecnici per ogni capacità?
\item[$\square$]
  Sono state inventariate route, job, consumer, webhook e strumenti
  amministrativi?
\item[$\square$]
  Sono noti database, writer, volumi, vincoli e consumer dello schema?
\item[$\square$]
  Sono state elencate integrazioni, librerie, licenze e dipendenze non
  supportate?
\item[$\square$]
  Esistono baseline di latenza, throughput, errori, disponibilità e
  costi?
\item[$\square$]
  Sono documentati RTO, RPO, retention e requisiti di sicurezza?
\item[$\square$]
  I comportamenti impliciti e i casi limite sono stati verificati con
  client reali o test?
\item[$\square$]
  Le incertezze sono attività con owner, evidenza attesa e scadenza?
\end{itemize}

\subsection{5.17.2 Checklist di
progettazione}\label{checklist-di-progettazione}

\begin{itemize}
\tightlist
\item[$\square$]
  La strategia (big bang, componenti, strangler o ibrida) è motivata dal
  rischio?
\item[$\square$]
  Ogni unità migrabile ha un confine e un proprietario dei dati?
\item[$\square$]
  È definito come instradare traffico e come disattivare il nuovo
  percorso?
\item[$\square$]
  I contratti API/dati hanno versioni, esempi e test?
\item[$\square$]
  Il modello di autenticazione e la migrazione delle sessioni sono
  approvati?
\item[$\square$]
  Le modifiche schema sono compatibili e il rollback è possibile?
\item[$\square$]
  Sono stati definiti timeout, retry, idempotenza e cancellazione?
\item[$\square$]
  Le differenze intenzionali rispetto alla sorgente sono state
  approvate?
\end{itemize}

\subsection{5.17.3 Checklist prima del
cutover}\label{checklist-prima-del-cutover}

\begin{itemize}
\tightlist
\item[$\square$]
  Build, test, scansioni e pubblicazione passano da checkout pulito.
\item[$\square$]
  Il target è stato avviato in un ambiente simile alla produzione.
\item[$\square$]
  Test di caratterizzazione e differenziali non hanno differenze
  critiche.
\item[$\square$]
  Sono verificati status, JSON, date, autorizzazioni, paginazione e
  limiti.
\item[$\square$]
  Dashboard, alert, log strutturati e correlation id funzionano.
\item[$\square$]
  Health check e shutdown sono provati con il load balancer.
\item[$\square$]
  Backup e restore sono stati testati; lo schema è compatibile con
  rollback.
\item[$\square$]
  Runbook identifica responsabili, soglie, canali e comandi di ritorno.
\item[$\square$]
  Il canary ha una durata, percentuale e criterio per proseguire.
\item[$\square$]
  Le operazioni di scrittura non vengono replicate senza idempotenza.
\end{itemize}

\subsection{5.17.4 Checklist per la
dismissione}\label{checklist-per-la-dismissione}

\begin{itemize}
\tightlist
\item[$\square$]
  Non esiste traffico, job, webhook o client che dipenda ancora dalla
  sorgente.
\item[$\square$]
  I dati sono riconciliati e la retention dei backup è approvata.
\item[$\square$]
  Credenziali, chiavi e route obsolete sono revocate.
\item[$\square$]
  Le dipendenze residue e la coesistenza hanno owner o data di
  rimozione.
\item[$\square$]
  Runbook, supporto e dashboard sono aggiornati sul sistema definitivo.
\item[$\square$]
  Costi e risorse del sistema precedente sono stati rimossi o motivati.
\end{itemize}

\subsection{5.17.5 Glossario essenziale}\label{glossario-essenziale}

{\def\LTcaptype{none} % do not increment counter
\begin{longtable}[]{@{}
  >{\raggedright\arraybackslash}p{(\linewidth - 2\tabcolsep) * \real{0.5000}}
  >{\raggedright\arraybackslash}p{(\linewidth - 2\tabcolsep) * \real{0.5000}}@{}}
\toprule\noalign{}
\begin{minipage}[b]{\linewidth}\raggedright
Termine
\end{minipage} & \begin{minipage}[b]{\linewidth}\raggedright
Significato nel capitolo
\end{minipage} \\
\midrule\noalign{}
\endhead
\bottomrule\noalign{}
\endlastfoot
Assessment & Raccolta di evidenze sul sistema, i suoi contratti, vincoli
e rischi \\
Baseline & Misure del sistema sorgente usate per confrontare il
target \\
Cutover & Passaggio del traffico o della proprietà di una capacità al
target \\
Strangler & Strategia che sostituisce un sistema per capacità,
mantenendo temporaneamente l'originale \\
Canary & Rilascio progressivo a una parte del traffico prima
dell'estensione \\
Parità funzionale & Equivalenza del comportamento osservabile per i
flussi concordati \\
Idempotenza & Proprietà per cui ripetere un'operazione non duplica
l'effetto \\
Outbox & Registro transazionale degli eventi da consegnare dopo il
commit \\
RTO/RPO & Obiettivo di tempo di ripristino e perdita dati massima
accettabile \\
Expand--migrate--contract & Evoluzione graduale di schema e codice
mantenendo compatibilità temporanea \\
\end{longtable}
}

\subsection{5.17.6 Gli errori da evitare}\label{gli-errori-da-evitare}

\begin{enumerate}
\def\labelenumi{\arabic{enumi}.}
\tightlist
\item
  \textbf{Riscrivere senza caratterizzare.} Il team si affida a ricordi
  e documenti incompleti; i test sul comportamento esistente vanno
  creati prima del cutover.
\item
  \textbf{Confondere parità e miglioramento.} Cambiare contratto, dati e
  architettura nello stesso rilascio rende impossibile isolare le cause
  delle regressioni.
\item
  \textbf{Tradurre i nomi invece delle responsabilità.} Un middleware o
  un service container con nome simile può avere lifecycle diverso;
  verificare semantica e scope.
\item
  \textbf{Ignorare il database condiviso.} Le tabelle possono essere un
  contratto di altri servizi; identificare tutti i writer e reader prima
  di migrare.
\item
  \textbf{Duplicare scritture senza protocollo.} Due sistemi non sono
  sincronizzati perché entrambi hanno ricevuto una chiamata; progettare
  idempotenza e riconciliazione.
\item
  \textbf{Confondere rollback applicativo e rollback dati.} La versione
  precedente può non interpretare i dati scritti dalla nuova.
\item
  \textbf{Ottimizzare prima di misurare.} AOT, cache o parallelismo
  aumentano complessità se il collo di bottiglia è altrove.
\item
  \textbf{Considerare test unitari come prova di produzione.} Database,
  proxy, identity e container richiedono test d'integrazione mirati.
\item
  \textbf{Lasciare la coesistenza senza data di uscita.} Il doppio
  sistema raddoppia costi e superficie di rischio; assegnare owner e
  milestone di dismissione.
\item
  \textbf{Fare del cutover una scommessa.} Soglie, responsabilità e
  rollback devono essere definiti e provati prima che il traffico venga
  spostato.
\end{enumerate}

\section{5.18 Appendice: piano di migrazione
riutilizzabile}\label{appendice-piano-di-migrazione-riutilizzabile}

Il seguente schema può essere copiato nel backlog di un progetto e
adattato. Ogni riga deve avere owner, data, evidenza e stato; non
lasciare caselle spuntate sulla sola base di una dichiarazione.

{\def\LTcaptype{none} % do not increment counter
\begin{longtable}[]{@{}
  >{\raggedright\arraybackslash}p{(\linewidth - 4\tabcolsep) * \real{0.3333}}
  >{\raggedright\arraybackslash}p{(\linewidth - 4\tabcolsep) * \real{0.3333}}
  >{\raggedright\arraybackslash}p{(\linewidth - 4\tabcolsep) * \real{0.3333}}@{}}
\toprule\noalign{}
\begin{minipage}[b]{\linewidth}\raggedright
Fase
\end{minipage} & \begin{minipage}[b]{\linewidth}\raggedright
Deliverable
\end{minipage} & \begin{minipage}[b]{\linewidth}\raggedright
Criterio di completamento
\end{minipage} \\
\midrule\noalign{}
\endhead
\bottomrule\noalign{}
\endlastfoot
Discovery & Inventario, baseline, rischi e contratti & Owner e
dipendenze identificati; lacune tracciate \\
Foundation & Repo, SDK, CI, ambiente, logging & Build pulita, deploy e
health check provati \\
Caratterizzazione & Fixture e test sorgente & Flussi critici
riproducibili e approvati \\
Vertical slice & Endpoint target completo & HTTP, dominio, storage e
test passano \\
Parità & Confronto sorgente/target & Differenze critiche risolte o
approvate \\
Canary & Traffico controllato & Soglie rispettate per la finestra
concordata \\
Cutover & Ownership e routing aggiornati & Target proprietario; rollback
ancora possibile \\
Stabilizzazione & Runbook e supporto & Incidenti e metriche entro
obiettivi \\
Dismissione & Sistema sorgente spento & Traffico nullo, accessi
revocati, dati conservati \\
\end{longtable}
}

Per ogni capacità, annotare un punteggio qualitativo di rischio
(basso/medio/alto) e le evidenze: non sommare numeri senza significato.
Se il criterio non è raggiunto, la decisione corretta può essere
fermarsi e ridurre l'incertezza, non accelerare il rilascio.

\backmatter
\end{document}
